% concatenated from Note0311_v0.tex
\documentclass[11pt]{amsart}
\usepackage{amsmath,amssymb,amsthm,geometry}
%\geometry{margin=1in}

\usepackage{latexsym, amsmath,amssymb}
\usepackage{graphicx}
\usepackage{hyperref}
\usepackage{enumerate}
\usepackage{color}
%\usepackage{showkeys}
%\newcommand{\cal}{\mathcal}
%\usepackage{bbm}
\newcommand{\newsection}[1]
{\subsection{#1}\setcounter{theorem}{0} \setcounter{equation}{0}
	\par\noindent}

%\setlength\evensidemargin{.5in}
%\setlength\textheight{44cc} \setlength\textwidth{30cc}
%\setlength\topmargin{0in} \setlength\parskip{5pt}

%\headheight=8pt     \topmargin=0pt
%\textheight=624pt   \textwidth=432pt
%\oddsidemargin=18pt \evensidemargin=18pt
%\usepackage[T1]{fontenc}
\usepackage[utf8]{inputenc}
\usepackage{amsmath,amssymb,amsthm,mathtools}
%\usepackage[margin=1in]{geometry}
\usepackage{microtype}

%\newtheorem{theorem}{Theorem}
%\newtheorem{lemma}{Lemma}
%\newtheorem{remark}{Remark}

%\newcommand{\R}{\mathbb{R}}
\newcommand{\norm}[2]{\lVert #1\rVert_{#2}}
\newcommand{\Lk}[1]{L_{#1}}
\newcommand{\Ek}[1]{E_{#1}}
\renewcommand{\epsilon}{\varepsilon}
\newtheorem{theorem}{Theorem}[section]
\newtheorem{lemma}[theorem]{Lemma}
\newtheorem{conj}[theorem]{Conjecture}
\newtheorem{definition}[theorem]{Definition}
\newtheorem{corr}[theorem]{Corollary}
\newtheorem{prop}[theorem]{Prop}
\newtheorem{proposition}[theorem]{Proposition}
\newtheorem{deff}[theorem]{Definition}
\newtheorem{remark}[theorem]{Remark}
\newcommand{\bth}{\begin{theorem}}
	\newcommand{\ble}{\begin{lemma}}
		\newcommand{\bcor}{\begin{corr}}
			\newcommand{\ltrt}{{L^2({\Bbb R}^3)}}
			\newcommand{\ltoo}{{L^2({\mathbb{R}}^3\backslash\mathcal{K})}}
			\newcommand{\bdeff}{\begin{deff}}
				\newcommand{\lirt}{{L^\infty({\Bbb R}^3)}}
				\newcommand{\lioo}{{L^{\infty}({\mathbb{R}}^3 \backslash \mathcal{K})}}
				\newcommand{\bprop}{\begin{proposition}}
					\newcommand{\ele}{\end{lemma}}
				\newcommand{\ecor}{\end{corr}}
			\newcommand{\edeff}{\end{deff}}
		\newcommand{\ii}{i}
		\newcommand{\eprop}{\end{proposition}}
	\newcommand{\rlnu}{{R_{\lambda}^{\nu}}}
	\newcommand{\trlnu}{{\tilde{R}_{\lambda}^{\nu}}}
	\newcommand{\tlnu}{{T_{\lambda}^{\nu}}}
	\newcommand{\ttlnu}{{\tilde{T}_{\lambda}^{\nu}}}
	\newcommand{\slnu}{{S_{\lambda}^{\nu}}}
	\newcommand{\slnut}{{}^t\!\slnu}
	\newcommand{\cd}{\, \cdot\, }
	\newcommand{\mlnu}{{m_{\lambda}^{\nu}}}
	\newcommand{\psilnu}{\psi_{\lambda}^{\nu}}
	\newcommand{\xilnu}{\xi_{\lambda}^{\nu}}
	\newcommand{\nlnu}{N_{\lambda}^{\nu}}
	\newcommand{\nl}{N_{\lambda}}
	\newcommand{\Rn}{{\mathbb R}^n}
	\newcommand{\jump}{{}}
	\newcommand{\la}{\lambda}
	\newcommand{\st}{{\Bbb R}^{1+3}_+}
	\newcommand{\eps}{\varepsilon}
	\newcommand{\e}{\varepsilon}
	\renewcommand{\l}{\lambda}
	\newcommand{\loc}{{\text{\rm loc}}}
	\newcommand{\comp}{{\text{\rm comp}}}
	\newcommand{\Coi}{C^\infty_0}
	\newcommand{\supp}{\text{supp }}
	\renewcommand{\Pi}{\varPi}
	\renewcommand{\Re}{ {\rm Re}\,}
	\renewcommand{\Im}{{\rm Im} \,}
	\renewcommand{\epsilon}{\varepsilon}
	\newcommand{\sgn}{{\text {sgn}}}
	\newcommand{\Gmid}{\Gamma_{\text{mid}}}
	\newcommand{\R}{\mathbb{R}}
	\newcommand{\N}{\mathbb{N}}
	\newcommand{\Rt}{{\Bbb R}^3}
	\newcommand{\Mdel}{{{\cal M}}^\alpha}
	\newcommand{\dist}{{{\rm dist}}}
	\newcommand{\Adel}{{{\cal A}}_\delta}
	\newcommand{\Kob}{{\cal K}}
	\newcommand{\Dia}{\overline{\Bbb E}^{1+3}}
	\newcommand{\Diap}{\overline{\Bbb E}^{1+3}_+}
	\newcommand{\Cyl}{{\Bbb E}^{1+3}}
	\newcommand{\Cylp}{{\Bbb E}^{1+3}_+}
	\newcommand{\Penrose}{{\cal P}}
	\newcommand{\Rplus}{{\Bbb R}_+}
	\newcommand{\parital}{\partial}
	\newcommand{\tidle}{\tilde}
	\newcommand{\grad}{\text{grad}\,}
	\newcommand{\ob}{{\mathcal K}}
%	\newcommand{\R}{{\mathbb R}}
	\newcommand{\ls}{\lesssim}
	\newcommand{\gs}{\gtrsim}
	\newcommand{\1}{{\rm 1\hspace*{-0.4ex}%
			\rule{0.1ex}{1.52ex}\hspace*{0.2ex}}}
	\newcommand{\CM}{\mathcal{C}_M}
	\newcommand{\ola}{\1_\la}
	
	\newcommand{\ala}{\tilde \1_\la}
	\newcommand{\Log}{{\rm Log}}
	%\newcommand{\square}{\Box}
	
	\newcommand{\subheading}[1]{{\bf #1}}
	\numberwithin{equation}{section}
		\newcommand{\logit}[1]{\log^{(#1)}}
\begin{document}

\title{Maximal growth of the Stein-Wainger oscillatory integral}
\author{ Cheng Zhang and Zhifei Zhu}

\address{Mathematical Sciences Center\\
	Tsinghua University\\
	Beijing, BJ 100084, China}
\email{czhang98@tsinghua.edu.cn; zhifeizhu@tsinghua.edu.cn}
\date{}
	\keywords{oscillatory integral; Gevrey class; Denjoy-Carleman}

\begin{abstract}
	We establish a precise hierarchy for the maximal growth of the Stein-Wainger oscillatory integral as the regularity of the phase varies over Denjoy-Carleman classes, such as the Gevrey classes and their generalizations. In particular, we resolve a problem posed by Wang--Zhang \cite{wz21}, motivated by eigenfunction restriction estimates on curves, and also provide a new proof of a theorem of Nagel--Wainger \cite{nw} on the Hilbert transform along curves. A key ingredient is the sharp estimate on the growth of a phase near a flat point. 
\end{abstract}
\maketitle
\section{Introduction}

Let $I=[-1,1]$, and let $\psi\in C^1(I)$ be a real-valued function. For $\lambda\gg1$, we define the Stein--Wainger oscillatory integral by
\begin{equation}\label{mla}
	m(\lambda)
	:=
	\mathrm{p.v.}\int_I e^{i\lambda\psi(t)}\,\frac{dt}{t}.
\end{equation}
The main objective of this paper is to identify the exact relationship between the regularity of $\psi$ and the maximal possible growth of $|m(\lambda)|$ as $\lambda\to\infty$. Throughout this paper, we assume that $\lambda$ is sufficiently large. We also note that, it is harmless to replace the interval $I$ by any bounded interval with $0$ in its interior, since the contribution of the integral outside a fixed neighborhood of $0$ is $O(1)$ and is therefore negligible for our purposes.

\subsection{Main results}
Our main results establish a new hierarchy of upper bounds, together with matching lower bounds, reflecting the regularity of the phase.

First, for $C^\alpha$ phases, the maximal growth of $m(\la)$ is logarithmic. 

\begin{proposition}\label{prop0}
	If $\psi\in C^\alpha(I)$ for some integer $\alpha\ge1$, then $m(\la)=O(\log\la)$. This bound is sharp: there exists a real-valued $\psi\in C^\alpha(I)$ such that $|m(\la)|\approx \log\la$.
\end{proposition}

Second, for $C^\infty$ phases, we may slightly improve the maximal growth to sub-logarithmic.
\begin{theorem}\label{thm:main0}
	If $\psi\in C^\infty(I)$, then $m(\la)=o(\log\la)$. This bound is sharp: for every integer $k\ge 2$, there exists a real-valued $\psi\in C^\infty(I)$ such that
	\[
	|m(\lambda_n)|\approx \frac{\log \lambda_n}{\logit{k}\lambda_n}	\quad\text{along a sequence }\lambda_n\to\infty.
	\]
\end{theorem}

Third, we may further improve the maximal growth for phases in the Gevrey classes. Recall that for $s\ge1$, a function $\psi\in C^\infty(I)$ belongs to the Gevrey class
$G^s(I)$ if there exists a constant $K>0$ such that for all $n\ge0$
\[
\sup_{t\in I} |\psi^{(n)}(t)|
\le K^{n+1} (n!)^s.
\]
In particular, $G^1=C^\omega$ is exactly the class of real-analytic functions. Thus Gevrey classes describe a regularity intermediate between smoothness and analyticity. Since their introduction in \cite{gev} in connection with regularity properties of the fundamental solution of the heat operator, Gevrey classes have been employed in a variety of contexts within the general theory of linear partial differential operators, including hypoellipticity, local solvability, and the propagation of singularities. For a comprehensive definition and a detailed discussion of Gevrey classes and their applications to linear partial differential operators, we refer the reader to \cite{rod93}.
\begin{theorem}\label{thm:main1}
	Let $\psi\in G^s(I)$ with $s>1$. Then $m(\la)=O(\log\log\la)$. 
	This bound is sharp: there exists a real-valued $\psi\in G^s(I)$ such that
	\[
	|m(\lambda_n)| \approx \log\log \lambda_n
	\quad\text{along a sequence }\lambda_n\to\infty.
	\]
\end{theorem}
For $\psi\in G^1(I)$, i.e. the analytic class, it is known that $m(\la)=O(1)$. See Pan \cite{pan93}.

Next, we introduce some refinements between the Gevrey classes and the analytic class. These allow us to see clearly how the maximal growth of $m(\la)$ improves from $O(\log\log\la)$ to $O(1)$. For $s\ge 1$ and an integer $k\ge 1$, a function $\psi\in C^\infty(I)$ belongs to the refined Gevrey class
$G_k^s(I)$ if there exists a constant $K>0$ such that for all sufficiently large $n$
\[
\sup_{t\in I} |\psi^{(n)}(t)|
\le K^{n+1} n!Q_{k,s}(n)^{n},
\]
where
\[
Q_{k,s}(n)=(\log^{(k)} n)^s\prod_{j=1}^{k-1}\log^{(j)} n
\]
and $\log^{(k)} n$ is the $k$-fold iterated logarithm. For instance, we have $Q_{1,s}(n)=(\log n)^s$ and $Q_{2,s}(n)=(\log n)(\log\log n)^s$. Thus these classes are refinements between $\bigcap_{s>1}G^{s}$ and $G^1=C^\omega$. Moreover, by the Denjoy--Carleman theorem (see H\"ormander \cite[Theorem 1.3.8]{hor1}), $G_k^s$ is quasianalytic if and only if $s=1$, since

\[
\sum_{n\gg 1}\frac{1}{n(\log n)\cdots(\log^{(k-1)}n)(\log^{(k)} n)^s}=\infty \iff s=1.
\]

The following theorem gives further improvements for the refined Gevrey classes and fills the gap between the Gevrey classes and the analytic class.

\begin{theorem}\label{thm:main2}
	 Let $\psi\in G_k^s(I)$ with $k\ge1$ and $s>1$. Then $m(\lambda)=O(\log^{(k+2)}\lambda).$
	This bound is sharp: there exists a real-valued $\psi\in G_k^s(I)$ such that
	\[
	|m(\lambda_n)| \approx \log^{(k+2)} \lambda_n
	\quad\text{along a sequence }\lambda_n\to\infty.
	\]
\end{theorem}

 If $\psi\in G_k^1(I)$, then $\psi$ is quasianalytic, and it is also known that $m(\lambda)=O(1)$ by Pan \cite{pan93}. Furthermore, if $P(t)$ is a real polynomial of degree $d$, then a much stronger result holds:
\begin{equation}\label{pbd}
	\Big|\mathrm{p.v.}\int_{-1}^1\frac{e^{i P(t)}}{t}dt\Big|\le C_d
\end{equation}
 where $C_d$ depends only on $d$, and the best constant satisfies $C_d\approx \log d$. This follows from the proof of Parissis \cite[(3.7), (4.4)]{pa08}. Similarly, if $R(t):=P(t)/Q(t)$ is a rational function where $P,Q$ are real polynomials, then 
 \begin{equation}\label{rbd}
 	\Big|\mathrm{p.v.}\int_{-1}^1\frac{e^{i R(t)}}{t}dt\Big|\le A
 \end{equation}
where $A$ depends only on the degrees of $P$ and $Q$ and not on their coefficients. This follows directly from the proof of Folch--Gabayet and Wright \cite[Theorem 1.1]{fw03}. Their results are about the integral on the whole real line, while the proof still works for bounded intervals.

\begin{remark}\label{rm1}{\rm The parameter $\alpha$ in Proposition \ref{prop0} and the parameter $s$ in Theorems \ref{thm:main1} and \ref{thm:main2} affect only the implicit constants and do not change the maximal growth in an essential way. Suppose that $\varphi$ is a smooth function on $\R^n$ and that $K$ is an odd, $-n$-homogeneous function on $\R^n$ that is integrable on the unit sphere. We have the standard reduction
	\begin{equation}
		\Big|\mathrm{p.v.}\int_{|x|\le 1}e^{i\la\varphi(x)}K(x)dx\Big|\le \frac12\|K\|_{L^1(S^{n-1})}\sup_{\theta\in S^{n-1}}\Big|\mathrm{p.v.}\int_{-1}^1\frac{e^{i\la\varphi(t\theta )}}tdt\Big|.
	\end{equation}
	Thus higher-dimensional analogues also hold whenever the derivatives of the phase $\varphi(t\theta)$ with respect to $t$ satisfy estimates uniformly in $\theta\in S^{n-1}$. For instance, these hold when $\varphi$ is a real polynomial or rational function by  \eqref{pbd} and \eqref{rbd}. }
\end{remark}

\begin{remark}\label{rm2}{\rm
	Given the gap between the Gevrey classes and the $C^\infty$ class, it is also important to examine classes of smooth functions that lie outside every Gevrey class. For instance, Jézéquel \cite{jez} established a trace formula conjectured by Dyatlov--Zworski \cite{dz2016} for Anosov flows in dynamical systems that holds for certain intermediate regularity classes. We will discuss them in Section \ref{sec7} and their extensions in Section \ref{sec8}. We will see how the maximal growth of $m(\la)$ increases to the universal bound $o(\log\la)$ of the $C^\infty$ class as the classes become larger.
	}
\end{remark}

\subsection{Background}
The problem is directly related to the singular oscillatory integral operators $T_\lambda$ of the form
\begin{equation}\label{Tlambda}T_\lambda f(t)={\rm p.v.}\int e^{i\lambda\phi(t,s)}(t-s)^{-1}a(t,s)f(s)ds,\end{equation}
where $\phi$ is smooth, $\lambda$ is real, and $a\in C_0^\infty(\mathbb{R}^2)$.
These operators and their higher-dimensional analogues have been studied by Stein--Wainger \cite{sw70,sw78}, Phong--Stein \cite{ps86}, Ricci--Stein \cite{rs87}, Pan \cite{pan91}, Seeger \cite{seeger94}, and Carbery--P\'erez \cite{cp99}. Phong--Stein \cite[p.~117]{ps86} showed that uniform $L^2(\mathbb{R})$ estimates for $T_\lambda$ imply the $L^2(\mathbb{R}^2)$ boundedness of the Hilbert transform $\mathcal{H}$ along variable curves: 
\begin{equation}\label{ht}
	\|\mathcal{H}\|_{L^2(\mathbb{R}^2)\to L^2(\mathbb{R}^2)}\le \sup_{\lambda\in \mathbb{R}}\|T_\lambda\|_{L^2(\mathbb{R})\to L^2(\mathbb{R})}.
\end{equation} Here $\mathcal{H}$ is defined initially on functions in $C_0^\infty(\mathbb{R}^2)$ by
\[\mathcal{H}f(x)=\eta(x){\rm p.v.}\int_{-\delta}^{\delta}f(x_1-t,x_2-\phi(x_1,x_1-t))\frac{dt}{t},\]
where $\eta\in C_0^\infty(\mathbb{R}^2)$ and $\delta>0$ is suitably small. Pan \cite{pan91} proved that $T_\lambda$ is uniformly bounded on $L^2(\mathbb{R})$ if one imposes a weak finite type condition. Later, Seeger \cite{seeger94}, Carbery--P\'erez \cite{cp99} investigated certain flat cases in which the finite-type condition fails.

In the translation-invariant case $\phi(t,s)=\psi(t-s)$, Nagel--Vance--Wainger--Weinberg \cite{nvww} proved necessary and sufficient conditions under the assumption that $\psi$ is even (or odd) and convex. Nagel--Wainger \cite[Theorem 4.1]{nw} constructed an odd smooth function $\psi$ on $[-1,1]$ such that the Hilbert transform $\mathcal{H}$ along the curve $(t,\psi(t))$ is unbounded on $L^2(\mathbb{R}^2)$. Indeed, the function $\psi$ is flat at $0$, so the convolution kernel
\begin{equation}\label{nwker}
	m(x, y)={\rm p . v .} \int_{-1}^{1} e^{i(x t+y \psi(t))} \frac{d t}{t}
\end{equation}
is unbounded on $\mathbb{R}^2$. Note that $m(0,\la)$ coincides with $m(\la)$ defined in \eqref{mla} with $I=[-1,1]$. Thus \cite[Theorem 4.1]{nw} is also implied by the unboundedness of $m(\la)$. Pan \cite{pan93} proved boundedness under the assumption that $\psi$ does
not vanish to infinite order at $0$. Motivated by the study of eigenfunction restriction estimates, Wang--Zhang \cite[Problem 4 in Section 4]{wz21} asked whether $m(\lambda)$ remains bounded for arbitrary smooth $\psi$. 

Stein--Wainger \cite{sw70} considered this type of singular oscillatory integrals with $I=\mathbb{R}$, and they showed that if $\psi(t)$ is a polynomial of degree $d$, then
\[\Big|{\rm p . v .} \int_{\mathbb{R}} e^{i\lambda\psi(t)} \frac{d t}{t}\Big|\le C_d\]
where the constant $C_d$ depends only on $d$. See also Parissis \cite{pa08}, Papadimitrakis--Parissis \cite{pp10}, and Al-Qassem, Cheng, and Pan \cite{acp14}. For rational phases $\psi$, uniform bounds still hold; see Folch--Gabayet and Wright \cite{fw03,fw032}, Wang--Wu \cite{ww23}, and Al-Qassem, Cheng, and Pan \cite{acp25}. However, this integral may diverge for many real-analytic phases, such as $\exp(t)$ and $\arctan(t)$. Nevertheless, our results show that if we replace $\mathbb{R}$ by a bounded interval, then an interesting new hierarchy emerges for the maximal growth of $m(\la)$ with respect to the regularity of the phase $\psi$. Moreover, our results give a negative answer to the problem posed by Wang--Zhang \cite[Problem 4 in Section 4]{wz21}. Furthermore, our results imply that \eqref{nwker} is unbounded on the $y$-axis, so in particular this gives a new proof of \cite[Theorem 4.1]{nw} by Nagel--Wainger.

\subsection{Proof Outline.} First, we use a van der Corput estimate to handle phases of finite type and reduce the problem to flat phases. Second, we use either Bang's lemma or the Taylor--Legendre method to prove sharp estimates on the growth of a phase near a flat point. Third, we use these sharp flat-point estimates to establish the oscillatory integral bounds. In particular, we obtain an abstract result (Theorem \ref{thm:main3}) for general Denjoy--Carleman classes and use it to establish Theorems \ref{thm:main1} and \ref{thm:main2}. Finally, we explicitly construct examples in the relevant classes to show that the estimates are sharp.

\subsection{Paper Structure.} In Section \ref{sec2}, we prove some key lemmas on the phase and Proposition \ref{prop0}. In Section \ref{sec3}, we prove Theorem \ref{thm:main0} for general smooth phases. In Section \ref{sec4}, we prove Theorem \ref{thm:main3} for phases in the Denjoy--Carleman class. In Section \ref{sec5}, we prove Theorem \ref{thm:main1} for phases in the Gevrey classes. In Section \ref{sec6}, we prove Theorem \ref{thm:main2} for phases in the refined Gevrey classes. In Section \ref{sec7}, we discuss the intermediate regularity classes of smooth functions introduced by Jézéquel \cite{jez}. In Section \ref{sec8}, we further discuss larger classes of smooth functions that lie outside every Gevrey class. In Section \ref{sec9}, we prove the derivative bounds for the sharpness examples.



\subsection{Notation.} Throughout this paper, $X\ls Y$ means $X\le CY$ for some positive constant $C$ independent of $\la$. If $X\ls Y$ and $Y\ls X$, we write $X\approx Y$.  If $X\ge CY$ for some sufficiently large constant $C>1$, we write $X\gg Y$. 

\subsection{Acknowledgments.}
The authors would like to thank Shaozhen Xu for helpful discussions and comments. The authors are supported in part by the National Key R\&D Program of China 2024YFA1015300. C.Z. is also supported in part by NSFC Grant 12371097. Z.Z. is also supported in part by NSFC Grant 12501065.
 
\section{Preliminaries}\label{sec2}
We rewrite the principal value integral as
\begin{equation}\label{eq:pv}
	m(\lambda)
	=
	\int_0^1 \frac{e^{i\lambda\psi(t)}-e^{i\lambda\psi(-t)}}{t}\,dt.
\end{equation}
Let
\begin{equation}\label{phidef}
	\phi(t)=\psi(t)-\psi(-t).
\end{equation}
	If $\phi\equiv 0$, then $\psi$ is even and $m(\lambda)=0$ identically. If $\phi(t)$ has a finite order of vanishing at $0$, then standard
oscillatory integral estimates imply that $m(\lambda)$ is bounded. See e.g. Pan \cite[Lemma 2.3]{pan93}. We give a direct proof for completeness.

\begin{lemma}\label{finitetype}
	Suppose that $\phi$ has a finite order of vanishing at 0. Then $m(\la)=O(1)$.
\end{lemma}
\begin{proof}
	Assume $\phi\not\equiv 0$ and let $k$ be the smallest odd integer with $\phi^{(k)}(0)\neq 0$.
	Then Taylor's theorem gives
	\[
	\phi(t)=c\,t^k+O(t^{k+2}),\qquad c=\frac{\phi^{(k)}(0)}{k!}\neq 0.
	\]
	Hence there exists $\delta\in(0,1)$ and constants $c_0,C_0>0$ such that for $0<t\le\delta$,
	\begin{equation}\label{eq:finite-type}
		c_0 t^k\le |\phi(t)|\le C_0 t^k,
		\qquad
		|\phi^{(k)}(t)|\ge c_0.
	\end{equation}
	
	Split
	\[
	m(\lambda)=\int_0^\delta \frac{e^{i\lambda\psi(t)}-e^{i\lambda\psi(-t)}}{t}\,dt
	+\int_\delta^1 \frac{e^{i\lambda\psi(t)}-e^{i\lambda\psi(-t)}}{t}\,dt.
	\]
	The second integral is $O(1)$, so it remains to bound the first one, denoted by $m_0(\la)$, uniformly in $\lambda$.
	
	Let $I_j=[2^{-j-1}\delta,2^{-j}\delta]$, $t_j=2^{-j}\delta$ ($j\ge 0$). Dyadically decompose the integral and
	rescale $t=t_js$, $s\in[1/2,1]$. We obtain
	\[
	m_0(\lambda)=\sum_{j\ge 0} J_j(\lambda),
	\qquad
	J_j(\lambda):=\int_{1/2}^1 \frac{e^{i\lambda\psi(t_js)}-e^{i\lambda\psi(-t_js)}}{s}\,ds.
	\]
	Define the scale parameter $\Lambda_j:=\lambda\,t_j^k=\lambda\,\delta^k\,2^{-jk}.$
We claim that \begin{equation}\label{clm}
	|J_j|\ls \min \{\Lambda_j,\ \Lambda_j^{-1/k}\}.
\end{equation} We postpone the proof of this claim and use it to obtain a uniform bound for $m_0(\lambda)$.

 Let $j_*$ be such that $\Lambda_{j_*}\ge 1>\Lambda_{j_*+1}$, i.e.
$2^{j_*}\approx (\lambda\delta^k)^{1/k}$. Then 
\[
\sum_{j\le j_*}|J_j(\lambda)|
\lesssim \sum_{j\le j_*}\Lambda_j^{-1/k}
=\sum_{j\le j_*}(\lambda\delta^k)^{-1/k}2^{j}
\lesssim (\lambda\delta^k)^{-1/k}2^{j_*}
\lesssim 1,
\]
and
\[
\sum_{j> j_*}|J_j(\lambda)|
\lesssim \sum_{j>j_*}\Lambda_j
=\sum_{j>j_*}\lambda\delta^k\,2^{-jk}
\lesssim \lambda\delta^k\,2^{-k(j_*+1)}
\lesssim 1.
\]

Now it remains to prove the claim \eqref{clm}. On the one hand, using $|e^{ix}-e^{iy}|\le |x-y|$ and \eqref{eq:finite-type}, we obtain
	\[
	\big|e^{i\lambda\psi(t)}-e^{i\lambda\psi(-t)}\big|
	\le \la\,|\psi(t)-\psi(-t)|
	=\la\,|\phi(t)|
	\lesssim \la\,t^k \quad (0<t\le\delta).
	\]
 Hence for $s\in[1/2,1]$,
	\[
	\big|e^{i\lambda\psi(t_js)}-e^{i\lambda\psi(-t_js)}\big|
	\lesssim \la\,(t_js)^k\lesssim \la\,t_j^k=\Lambda_j.
	\]
	Therefore
	\begin{equation}\label{eq:small}
		|J_j(\lambda)|
		\ls \int_{1/2}^1 \frac{\Lambda_j}{s}\,ds
		\lesssim \Lambda_j.
	\end{equation}
	
	
	On the other hand, write $J_j=A_j-B_j$ where
	\[
	A_j(\lambda):=\int_{1/2}^1 \frac{e^{i\lambda\psi(t_js)}}{s}\,ds,
	\qquad
	B_j(\lambda):=\int_{1/2}^1 \frac{e^{i\lambda\psi(-t_js)}}{s}\,ds.
	\]
Let $\Phi_j(s)=\psi(t_js)$. 
	Since $k$ is odd and
\[
\phi^{(k)}(t)=\psi^{(k)}(t)+\psi^{(k)}(-t)=2\psi^{(k)}(0)+O(t),
\]
by \eqref{eq:finite-type} we have $|\psi^{(k)}(u)|\gtrsim 1$ for $|u|\le\delta$
	(up to adjusting $\delta$), hence $|\Phi_j^{(k)}(s)|\gtrsim t_j^k$ on $[1/2,1]$. Then a standard van der Corput estimate for phases with a nonvanishing $k$th derivative on a fixed interval yields
	\[
	|A_j(\lambda)|\lesssim (\lambda t_j^k)^{-1/k}=\Lambda_j^{-1/k}.
	\]
	Similarly $|B_j(\lambda)|\lesssim \Lambda_j^{-1/k}$, hence
	\begin{equation}\label{eq:large}
		|J_j(\lambda)|\le |A_j(\lambda)|+|B_j(\lambda)|\lesssim \Lambda_j^{-1/k}.
	\end{equation}
	This proves the claim \eqref{clm} and completes the proof.\end{proof}
	
	
	Next, we need the following one-dimensional Bang's lemma (Bang \cite[Equation (14)]{bang}, Bruna \cite[Lemma 2.5]{bruna}), which allows us to control the rate of increase of a function near a flat point. For completeness, we give a direct proof of the version we use.
\begin{lemma}[Bang]\label{lem:BB}
	Let $r>0$ and let $g\in C^\infty((-r,r))$ satisfy $g(x)=0\ (x\le 0).$
	Let $(A_n)_{n\ge 0}$ be a sequence of positive numbers such that
	\[
	\sup_{x\in (-r,r)}|g^{(n)}(x)|\le A_n\qquad (n\ge 0),
	\]
	and such that the quotient sequence
$\eta_n:=A_n/A_{n-1}\ (n\ge 1)$
	is nondecreasing. Fix $x\in(0,r)$ and an integer $\ell\ge 1$. If
	\[
	\sum_{j=\ell}^{\infty}\frac1{\eta_j}> 4x,
	\]
	then
	\[
	|g(x)|\le A_0 2^{-\ell}.
	\]
\end{lemma}

\begin{proof}The idea is to connect $x$ to $0$ by a chain of short intervals whose lengths are
	comparable to $\eta_j^{-1}$ and then to iterate the fundamental theorem of calculus from high derivatives down to lower ones. The monotonicity of ($\eta_j$) is exactly what allows each step to close.
	
	
	Choose $n\ge \ell$ so that
	\[
	\sum_{j=\ell}^n \frac1{\eta_j}\ge 4x.
	\]
	Set
	\[
	a:=\frac{x}{\sum_{j=\ell}^n \eta_j^{-1}}\le \frac14.
	\]
	Choose points
	\[
	t_\ell=x>t_{\ell+1}>\cdots > t_{n+1}=0
	\]
	so that
	\[
	t_j-t_{j+1}=\frac{a}{\eta_j}\qquad (\ell\le j\le n).
	\]
	For $\ell\le j\le n$ set $I_j=[t_{j+1},t_j]$ and let $I_{n+1}=(-r,0]$. For integers
	$0\le p\le q\le n+1$, define
	\[
	F(p,q):=\sup_{x\in I_q}|g^{(p)}(x)|.
	\]
	We claim that
	\begin{equation}\label{eq:claim}
		F(p,q)\le (2a)^{q-p}A_p
		\qquad (\ell\le q\le n+1,\; 0\le p\le q).
	\end{equation}
	This is clear when $p=q$, and also when $q=n+1$, because $g$ and all its derivatives vanish on
	$I_{n+1}=(-r,0]$.
	
	Now let $\ell\le q\le n$ and $0\le p<q$. For every $y\in I_q$, the fundamental theorem of calculus gives
	\[
	|g^{(p)}(y)|
	\le |g^{(p)}(t_{q+1})| + (t_q-t_{q+1})\sup_{x\in I_q}|g^{(p+1)}(x)|.
	\]
	Hence
	\[
	F(p,q)\le F(p,q+1)+\frac{a}{\eta_q}F(p+1,q).
	\]
	Assuming \eqref{eq:claim} already known for $(p,q+1)$ and $(p+1,q)$, we get
	\[
	F(p,q)
	\le (2a)^{q+1-p}A_p+\frac{a}{\eta_q}(2a)^{q-p-1}A_{p+1}.
	\]
	Because $(\eta_j)$ is nondecreasing and $p+1\le q$, we have $A_{p+1}=A_p\eta_{p+1}\le A_p\eta_q$, so
	\[
	F(p,q)
	\le (2a)^{q-p}A_p\Bigl(2a+\frac12\Bigr)
	\le (2a)^{q-p}A_p,
	\]
	since $a\le 1/4$. This proves \eqref{eq:claim}. Finally, since $x=t_\ell\in I_\ell$, \eqref{eq:claim} with
	$(p,q)=(0,\ell)$ gives
	\[|g(x)|\le F(0,\ell)\le A_0(2a)^\ell\le A_0 2^{-\ell}.\]
\end{proof}


	\subsection{Proof of Proposition \ref{prop0}}
	Suppose that $\psi\in C^\alpha([-1,1])$ for some integer $\alpha\ge1$. Using $|e^{ix}-e^{iy}|\le \min\{2,|x-y|\}$, \eqref{eq:pv}, and $|\phi(t)|\ls |t|$, we obtain
	\[
	|m(\lambda)|
	\ls
	\int_0^1 \frac{\min\{1,\la t\}}{t}\,dt\ls \log\la.
	\]
We remark that one may obtain a more precise estimate, namely $|m(\la)|\le \frac{2}{\alpha}\log\la+O(1)$, by using a refined argument similar to the proof of Theorem \ref{thm:main0}.

Next, we choose $\psi(t)=\max(t^{\alpha+1},0)\in C^{\alpha}([-1,1])$ and then
\begin{align*}
	m(\lambda)
	&=
	\int_0^1 \frac{e^{i\lambda t^{\alpha+1}}-1}{t}\,dt=\frac1{\alpha+1}\int_0^\la\frac{e^{iu}-1}{u}\,du\\
	&=
	-\frac1{\alpha+1}\log\lambda + O(1).
\end{align*}
This proves the sharpness of the upper bound. For large $\alpha$, this also shows that the coefficient $\approx\frac1\alpha$ is essentially sharp.
	\section{General smooth functions}\label{sec3}
	We prove Theorem \ref{thm:main0} in this section. Let $\psi\in C^\infty([-1,1])$ and $\phi(t)=\psi(t)-\psi(-t)$. By Lemma \ref{finitetype}, it remains to handle the case in which $\phi$ is flat at $0$.
	Then for every $N\ge1$ there exist constants $C_N,\delta_N>0$ such that
	\[
	|\phi(t)|\le C_N |t|^N
	\qquad (|t|\le\delta_N).
	\]
	Hence
	\[
	|e^{i\lambda\psi(t)}-e^{i\lambda\psi(-t)}|
	\le \lambda C_N |t|^N.
	\]
	
	Split \eqref{eq:pv} at
	\[
	r:=\min\bigl(\delta_N,(\lambda C_N)^{-1/N}\bigr).
	\]
	Then
	\[
	|m(\lambda)|
	\le \int_0^r \lambda C_N t^{N-1}\,dt + 2\int_r^1 \frac{dt}{t}
	\le \frac1N + \frac{2}{N}\log\lambda + O(1).
	\]
	Dividing by $\log\lambda$ and then letting $\lambda\to\infty$, followed by $N\to\infty$, yields $m(\la)=o(\log\la)$ as claimed.



\subsection{Sharpness}
We construct examples to show that the sub-logarithmic growth $o(\log\la)$ is sharp. The construction uses a ``dyadic plateau'' phase for which $m(\lambda)$
grows linearly in an index $n$ along a chosen frequency subsequence
$\lambda_n$, and then we tune the subsequence so that
$n\approx \log\lambda_n/\logit{k}\lambda_n$.



Fix $\eta\in C_c^\infty((1/2,1))$ with $0\le \eta\le 1$ and
\[
\eta\equiv 1\ \text{ on } [3/5,4/5].
\]
For $j\ge 1$ set
\[
\eta_j(t):=\eta(2^j t),
\qquad
J_j:=\Big[\frac{3}{5\cdot 2^j},\frac{4}{5\cdot 2^j}\Big].
\]
Then $\eta_j\equiv 1$ on $J_j$ and $\operatorname{supp}\eta_j\subset(2^{-(j+1)},2^{-j})$,
so the supports are pairwise disjoint.

Let $\{q_j\}_{j\ge 1}$ be an increasing sequence of odd integers with $q_j\ge 3$.
Define
\[
Q_j:=\prod_{\ell=1}^j q_\ell,
\qquad
a_j:=\frac{\pi}{Q_j}.
\]
Define the phase
\begin{equation}\label{eq:psi-def}
	\psi(t):=
	\begin{cases}
		0, & t\le 0,\\[4pt]
		\displaystyle \sum_{j=1}^\infty a_j\,\eta_j(t), & t>0.
	\end{cases}
\end{equation}
Because the supports of $\eta_j$ are disjoint, the sum is locally finite on $(0,1)$.

Now we choose $q_j$ to force a given iterated-log profile. For $k\ge 2$, fix a large integer $j_0$ so that $\logit{k-2}(j+j_0)\ge 2$ for all $j\ge 1$, and define
	\begin{equation}\label{eq:qj-choice}
		q_j:=2\Big\lfloor \logit{k-2}(j+j_0)\Big\rfloor+1.
	\end{equation}
	Then $q_j$ is odd, $q_j\ge 5$, and
$\log q_j \approx \logit{k-1} j$ and $\log Q_n \approx
	n\,\logit{k-1} n$. 

\begin{lemma}\label{lem:smooth}
	With the choice of $q_j$ above,
	the function $\psi$ defined by \eqref{eq:psi-def} lies in $C^\infty(\mathbb R)$.
\end{lemma}

\begin{proof}
	Fix an integer $m\ge 0$.
	On $\operatorname{supp}\eta_j$ we have $\psi(t)=a_j\eta_j(t)$ and
	\[
	\psi^{(m)}(t)=a_j\,2^{jm}\,\eta^{(m)}(2^j t).
	\]
Since $\log Q_j\approx j\log^{(k-1)}j$, we obtain
	\[
	\sup_{t\in(0,2^{-\ell})}|\psi^{(m)}(t)|
	\lesssim \sup_{j\ge \ell}\frac{2^{jm}}{Q_j}\to 0\quad (\ell \to\infty).
	\]
	Therefore $\psi^{(m)}(t)\to 0$ as $t\to 0^+$ for every $m$.
	Since $\psi\equiv 0$ on $(-\infty,0]$, all derivatives match at $0$ and
	$\psi\in C^\infty(\mathbb R)$.
\end{proof}

Let $\lambda_n=Q_n$. Since $\log \lambda_n\approx n\log^{(k-1)}n$, we have $n\approx \log \lambda_n/\log^{(k)}\lambda_n$ for all large $n$.

\begin{lemma}[Lower bound]\label{lem:lower}
	There exists $c_0>0$ such that for all $n\ge 1$,
	\[
	|m(\lambda_n)|\ge c_0\,n.
	\]
\end{lemma}

\begin{proof}
Since $\psi(-t)=0$ for $t>0$ by construction, we have
	\[
	-\Re m(\lambda)=\int_0^1 \frac{1-\cos(\lambda\psi(t))}{t}\,dt\ge 0.
	\]
	Fix $n$ and take $t\in J_j$ with $1\le j\le n$. Then $\eta_j(t)=1$ and
	$\eta_\ell(t)=0$ for $\ell\neq j$, so $\psi(t)=a_j=\pi/Q_j$.
	Therefore
	\[
	\lambda_n\psi(t)=Q_n\cdot\frac{\pi}{Q_j}
	=\pi\prod_{\ell=j+1}^n q_\ell.
	\]
	Since each $q_\ell$ is odd, we have
	$\cos(\lambda_n\psi(t))=-1$ on $J_j$.
	Thus
	\[
	-\Re m(\lambda_n)
	\ge \sum_{j=1}^n \int_{J_j}\frac{2}{t}\,dt
	=2\sum_{j=1}^n \log\!\left(\frac{ \frac{4}{5\cdot 2^j} }{ \frac{3}{5\cdot 2^j} }\right)
	=2n\log\!\left(\frac43\right).
	\]
	So $|m(\lambda_n)|\ge -\Re m(\lambda_n)\ge c_0 n$ with $c_0:=2\log(4/3)$.
\end{proof}

\begin{lemma}[Upper bound]\label{lem:upper}
	There exists $C_0>0$ such that for all $n\ge 1$,
	\[
	|m(\lambda_n)|\le C_0\,n.
	\]
\end{lemma}

\begin{proof}
	Write $m(\lambda_n)=\int_0^1 \frac{e^{i\lambda_n\psi(t)}-1}{t}\,dt$.
	Split $(0,1)$ into dyadic shells $I_j:=(2^{-(j+1)},2^{-j})$.
	On each $I_j$ we have $|e^{i\lambda_n\psi(t)}-1|\le 2$ and
	\[
	\int_{I_j}\frac{dt}{t}=\log 2,
	\]
	so
	\[
	\int_{I_j}\frac{|e^{i\lambda_n\psi(t)}-1|}{t}\,dt \le 2\log 2.
	\]
	Summing this over $j=1,\dots,n$ gives a bound $\lesssim n$.
	
	For the tail $j>n$, we use $|e^{ix}-1|\le |x|$ and the fact that on $I_j$
	the phase equals $a_j\eta_j$ (or $0$), hence $|\psi(t)|\le |a_j|=\pi/Q_j$.
	Thus for $j>n$,
	\[
	\int_{I_j}\frac{|e^{i\lambda_n\psi(t)}-1|}{t}\,dt
	\le \int_{I_j}\frac{|\lambda_n|\,|\psi(t)|}{t}\,dt
	\le \frac{\pi Q_n}{Q_j}\int_{I_j}\frac{dt}{t}
	= \pi\log 2\cdot \frac{Q_n}{Q_j}.
	\]
	Since $q_\ell\ge 3$, we have $Q_j\ge Q_n\,3^{j-n}$ for $j>n$, hence
	$\sum_{j>n}Q_n/Q_j\le \sum_{r\ge 1}3^{-r}\ls 1$.
	Therefore the tail contributes $O(1)$ uniformly in $n$.
	Combining, $|m(\lambda_n)|\le C_0 n$.
\end{proof}


Thus, by Lemmas~\ref{lem:lower} and \ref{lem:upper}, we obtain
\[
|m(\lambda_n)|\approx n
\approx \frac{\log \lambda_n}{\logit{k}\lambda_n}.
\]

\section{Denjoy--Carleman class}\label{sec4}
In this section, we prove general upper bounds for $m(\la)$ when the phase belongs to a Denjoy--Carleman class associated with a log-convex sequence.

Let $M=(M_n)_{n\ge 0}$ be a positive sequence with $M_0=1$. We say that $M$ is \emph{log-convex} if
\[
M_n^2\le M_{n-1}M_{n+1}\qquad (n\ge 1).
\]
For such a sequence, define the quotient sequence
\[
\mu_n:=\frac{M_n}{M_{n-1}}\qquad (n\ge 1).
\]
Log-convexity is equivalent to the monotonicity of $(\mu_n)_{n\ge 1}$.


For a log-convex sequence $M$, a function $\psi\in C^\infty([-1,1])$ belongs to the Denjoy--Carleman class
$\CM([-1,1])$ if there exists a constant $K>0$ such that for all $n\ge0$
\begin{equation}\label{defdc}
	\sup_{t\in[-1,1]}|\psi^{(n)}(t)|\le K^{n+1}M_n.
\end{equation}
The Gevrey class $G^s$ and its refinements $G_k^s$ are important special cases of the Denjoy--Carleman class. We first establish a general theorem for the Denjoy-Carleman class in this section, and then discuss concrete estimates and sharpness examples for the Gevrey class $G^s$ and its refinements $G_k^s$ in the next two sections.

The central quantity below is the tail function
\begin{equation}\label{eq:TM}
	T_M(N):=\sum_{j\ge N}\frac1{\mu_j},
\end{equation}
with the convention $T_M(N)=\infty$ if the series diverges. By the Denjoy--Carleman theorem (see H\"ormander \cite[Theorem 1.3.8]{hor1}), the class is quasianalytic exactly when $T_M(1)=\infty$. Specifically, for the Gevrey class $G^s$ with $s>1$, we have $T_M(N)\approx N^{1-s}$, and for the refined Gevrey class $G_k^s$ with $s>1$, we have $T_M(N)\approx (\log^{(k)}N)^{1-s}$.

\begin{theorem}\label{thm:main3}
	Let $M=(M_n)_{n\ge 0}$ be a log-convex sequence, and let $\mu_n=M_n/M_{n-1}$ and $T_M$ be as in \eqref{eq:TM}.
	
	\smallskip
	\noindent
	\textup{(i)} If $T_M(1)=\infty$, then $m(\lambda)=O(1).$
	
	\smallskip
	\noindent
	\textup{(ii)} Assume now that $T_M(1)<\infty$ and $\mu_NT_M(N)\ge C_0$ for some $C_0>1$ and all large $N$.
	Then there exist constants $C,c>0$ such that
	\begin{equation}\label{ub3}
		|m(\lambda)|\le C\Bigl(1+\log \frac1{T_M(c\log \lambda)}\Bigr)
		\qquad (\lambda\ge 2).
	\end{equation}
\end{theorem}
The condition ``$\mu_NT_M(N)\ge C_0$'' roughly means that the series $\sum_{j\ge N}1/\mu_j$ cannot converge faster than a geometric series. It can be relaxed, but it is enough for our purposes, since it holds for the Gevrey class $G^s$ $(s>1)$ and the refined Gevrey class $G_k^s$ $(s>1)$.

To prove the theorem, we need to control the rate of growth of $\phi$ near 0. We achieve this by using Bang's Lemma \ref{lem:BB}.
\begin{proposition}[Flat-point estimate]\label{prop:flat}
	Let $M=(M_n)_{n\ge 0}$ be a log-convex sequence. Let
	\[
	N_M(r):=\sup\bigl\{N\ge 1:\ T_M(N)\ge r\bigr\}\in \N\cup\{0,\infty\}.
	\]
	Let $\psi\in \CM([-1,1])$ and
	\[
	\phi(t):=\psi(t)-\psi(-t).
	\]
	Assume that $\phi$ is flat at $0$. Then there exist constants $C,c>0$ such that for all sufficiently small $|t|$
	\begin{equation}\label{eq:flat-bound}
		|\phi(t)|\le C\,2^{-N_M(c|t|)}.
	\end{equation}
	If in addition $T_M(1)=\infty$, then $\phi\equiv 0$ on $[-1,1]$.
\end{proposition}
Note that $N_M(r)$ is roughly the ``inverse'' of the function $r=T_M(N)$.
\begin{proof}
	Extend $\phi$ by zero to the left:
	\[
	F(x):=
	\begin{cases}
		0,& x\le 0,\\
		\phi(x),& 0<x<1.
	\end{cases}
	\]
	Since $\phi$ is flat at $0$, the extension $F$ belongs to $C^\infty([-1,1])$ and obeys
	\[
	\sup_{x\in[-1,1]}|F^{(n)}(x)|\le 2K^{n+1}M_n.
	\]
	Apply Lemma \ref{lem:BB} with $A_n:=2K^{n+1}M_n$. Then
	\[
	\eta_n=\frac{A_n}{A_{n-1}}=\frac{KM_n}{M_{n-1}}=K\mu_n,
	\]
	which is nondecreasing by log-convexity. Therefore, if $x\in(0,1)$ and $T_M(\ell)>4Kx$, then
	\[
	|F(x)|\le 2\cdot 2^{-\ell}=2^{1-\ell}.
	\]
	Equivalently, there are constants $C,c>0$ such that
	\[
	|\phi(x)|\le C\,2^{-N_M(cx)}
	\qquad (0<x\ll 1).
	\]
	Applying the same argument to $x\mapsto \phi(-x)$ gives \eqref{eq:flat-bound} for both signs of $t$.
	
	If $T_M(1)=\infty$, then $T_M(\ell)=\infty$ for every $\ell$. Fix $x\in(0,1)$ and let $\ell\to\infty$ in Lemma \ref{lem:BB}. We obtain $F(x)=0$. Thus $\phi(x)=0$ for $x>0$, and by symmetry also for $x<0$.
\end{proof}
\begin{remark}\label{rm3}{\rm It is natural to estimate the flat functions by Taylor's theorem and Legendre transform. Indeed, since $\phi$ is flat at 0, we have
\begin{equation}\label{taylor}
	|\phi(t)|\le \inf_{n\ge1}\frac{K^{n+1}M_n |t|^n}{n!}=K\exp\Big(-\Phi^*\Big(\log\frac1{K|t|}\Big)\Big)
\end{equation}
where $\Phi(n)=\log(M_n/n!)$ and 
\begin{equation}\label{leg}
	\Phi^*(y) := \sup_{n \geq 1}\{ny - \Phi(n)\}
\end{equation}
is the Legendre transform of $\Phi$. This Taylor--Legendre method works for the Gevrey classes and yields Proposition \ref{prop2}, but it cannot produce the correct estimates for the refined Gevrey classes in Proposition \ref{prop1}. For example, if $\phi\in G_2^2$, the Taylor--Legendre method only gives, for some constant $c>0$,
\[|\phi(t)|\ls \exp\Big(-\exp\Big(\frac{c}{|t|(\log\frac1{|t|})^2}\Big)\Big),\ \ |t|\ll1,\] 
while the correct estimate is given by Proposition \ref{prop1}: for some constants $B,\gamma>0$
\[|\phi(t)|\ls \exp\Big(-B\exp\Big(\exp\Big(\gamma |t|^{-1}\Big)\Big)\Big),\ \ |t|\ll1.\]
Nevertheless, we observe that this Taylor--Legendre method can still produce correct estimates for the intermediate regularity classes of smooth functions that lie outside every Gevrey class. See Remark \ref{rm4}.
}
\end{remark}


\subsection{Proof of Theorem \ref{thm:main3}}
	
By the Denjoy--Carleman theorem (see H\"ormander \cite[Theorem 1.3.8]{hor1}) and Lemma \ref{finitetype}, we only need to prove part (ii) of Theorem~\ref{thm:main3} when $\phi$ is flat at $0$. By Proposition~\ref{prop:flat}, after shrinking constants if needed, there exist $C_1,c_1>0$ such that
\[
|\phi(t)|\le C_1 2^{-N}
\qquad \text{whenever } 0<t\le t_N:=c_1 T_M(N).
\]
Since $T_M(N)\downarrow 0$, the intervals $I_N:=(t_{N+1},t_N]\ (N\gg 1)$
form a shell decomposition near the origin. We write
\[
|m(\lambda)|\lesssim 1+\sum_{N\ge N_0}\int_{I_N}\min\{2,\lambda|\phi(t)|\}\,\frac{dt}{t}
\]
for some fixed large $N_0$. On $I_N$ we have $|\phi(t)|\le C_1 2^{-N}$, hence
\begin{equation}\label{eq:shellsum}
	|m(\lambda)|
	\lesssim 1+\sum_{N\ge N_0}\min\{1,\lambda 2^{-N}\}\log\frac{T_M(N)}{T_M(N+1)}.
\end{equation}
For $N\ls\log\la$, we obtain for some $c>0$
\begin{align*}
	\sum_{N_0\le N\ls \log\la}\log\frac{T_M(N)}{T_M(N+1)}
	\lesssim 1+\log\frac1{T_M(c\log\la+1)}.
\end{align*}

For $N\gg \log\la$, 
by the assumption that $T_M(N)\ge C_0/\mu_N$ for some $C_0>1$ and all large $N$, we have for all large $N$
\[
T_M(N+1)=T_M(N)\Bigl(1-\frac1{\mu_NT_M(N)}\Bigr)\approx T_M(N).
\]
Hence
\[
\sum_{N\gg \log\la}\lambda 2^{-N}\log\frac{T_M(N)}{T_M(N+1)}
\lesssim \sum_{N\gg \log\la}\lambda2^{-N}
\lesssim 1.
\]
Combining the two ranges yields
\[
|m(\lambda)|\ls 1+\log\frac1{T_M(c\log\lambda)}.
\]


\section{Gevrey class}\label{sec5}

We prove Theorem \ref{thm:main1} in this section. For $s\ge1$, fix $\psi\in G^s([-1,1])$ and define $\phi(t)=\psi(t)-\psi(-t)$. By Lemma \ref{finitetype}, it remains to handle the case in which $\phi$ is flat at $0$.
We need to estimate the rate of growth of $\phi$ near $0$. Since $M_n=(n!)^s$, we can apply Proposition \ref{prop:flat} with 
$$T_M(N)\approx N^{1-s}\quad \text{and}\quad N_M(r)\approx r^{-1/(s-1)}$$ 
when $s>1$. Note that $T_M(1)=\infty$ when $s=1$.


\begin{proposition}[Flat-point estimate]\label{prop2}
	Let $s\ge1$, and let $\phi\in G^s([-1,1])$ be flat at $0$, i.e. $\phi^{(n)}(0)=0\ (n\ge 0).$
Then the following statements hold.

\smallskip
\noindent
\textup{(i)} If $s=1$, then $\phi\equiv 0$ on $[-1,1]$.

\smallskip
\noindent
\textup{(ii)} If $s>1$, then there exist constants $A,B>0$ such that for all sufficiently small $|t|$,
\begin{equation}\label{ub2}
	|\phi(t)|\le A\exp\bigl(-B |t|^{-1/(s-1)}\bigr).
\end{equation}
\end{proposition}
The estimate \eqref{ub2} is essentially sharp, since it is classical that $\exp(-|t|^{-1/(s-1)})$ is flat at $0$ and belongs to $G^s([-1,1])$. By Theorem \ref{thm:main3} we have $m(\la)=O(1)$ when $s=1$ and $m(\la)=O(\log\log\la)$ when $s>1$.
This proves the upper bound in Theorem~\ref{thm:main1}.

\subsection{ Sharpness}

Fix $s>1$ and set $a=\frac{1}{s-1}$.
Define
\[
\psi(t)=
\begin{cases}
	0,& t\le 0,\\[2pt]
	e^{-t^{-a}},& t>0.
\end{cases}
\]
It is classical that $\psi\in G^s([-1,1])$ and is flat at $0$.

For this choice,
\[
m(\lambda)
=
\int_0^1 \frac{e^{i\lambda\psi(t)}-1}{t}\,dt.
\]
With the substitution $u=\psi(t)=e^{-t^{-a}}$, one computes
\[
\frac{dt}{t}=(s-1)\frac{du}{u\log(1/u)},
\]
hence
\[
m(\lambda)
=
(s-1)\int_0^{e^{-1}}
\frac{e^{i\lambda u}-1}{u\log(1/u)}\,du.
\]
Let $\lambda_n=2\pi n$ and
\[
I_j:=\Bigl[\frac{j+1/4}{n},\,\frac{j+3/4}{n}\Bigr],\quad j=1,2,3,\dots.
\]
On \(I_j\), we have \(1-\cos(2\pi n u)\ge 1\).
Since the weight
\[
w(u):=
\frac1{u\,\log(1/u)}
\]
is decreasing for small \(u\), summing over \(1\le j\ls n\) yields, for some fixed $c_0\in (0,1/e)$,
\[
-\Re m(\lambda_n)
\gtrsim
(s-1)\int_{1/n}^{c_0} w(u)\,du\gs (s-1)\log\log n.
\]
Therefore
\[
|m(\lambda_n)|\ge -\Re m(\lambda_n) \gtrsim (s-1)\log\log \lambda_n.
\]
This shows that the upper bound $O(\log\log\la)$ is sharp.


\section{Refined Gevrey class}\label{sec6}
We prove Theorem \ref{thm:main2} in this section. For $s\ge1$ and an integer $k\ge1$, fix $\psi\in G_k^s([-1,1])$ and define $\phi(t)=\psi(t)-\psi(-t)$. By Lemma \ref{finitetype}, we only need to consider the case in which $\phi$ is flat at $0$. Since $M_n=(\log^{(k)} n)^s\prod_{j=1}^{k-1}\log^{(j)} n$ for large $n$, we can apply Proposition \ref{prop:flat} with 
$$T_M(N)\approx (\log^{(k)}N)^{1-s}\quad \text{and}\quad N_M(r)\approx \exp^{(k)}(r^{-1/(s-1)})$$ 
when $s>1$. Here $\exp^{(k)}$ is the $k$-fold iterated exponential. Note that $T_M(1)=\infty$ when $s=1$.

\begin{proposition}[Flat-point estimate]\label{prop1}
	For $s\ge1$ and integer $k\ge1$, let $\phi\in G_k^s([-1,1])$ be flat at $0$, i.e. $\phi^{(n)}(0)=0\ (n\ge 0).$
	Then the following statements hold.
	
	\smallskip
	\noindent
	\textup{(i)} If $s=1$, then $\phi\equiv 0$ on $[-1,1]$.
	
	\smallskip
	\noindent
	\textup{(ii)} If $s>1$, then there exist constants $A,B,\gamma>0$ such that for all sufficiently small $|t|$,
	\begin{equation}\label{ub1}
		|\phi(t)|\le A\exp\bigl(-B \exp^{(k)}(\gamma |t|^{-1/(s-1)})\bigr).
	\end{equation}
\end{proposition}
The estimate \eqref{ub1} is essentially sharp, since $\exp(-\exp^{(k)}(|t|^{-1/(s-1)}))$ is flat at $0$ and belongs to $G_k^s([-1,1])$ (see Section \ref{sec9.1}). By Theorem \ref{thm:main3} we have $m(\la)=O(1)$ when $s=1$ and $m(\la)=O(\log^{(k+2)}\la)$ when $s>1$.
This proves the upper bound in Theorem~\ref{thm:main2}.


\subsection{Sharpness}\label{6harp}
For convenience, we denote $E_1(x)=e^x,\ E_{j+1}(x)=\exp(E_j(x)),\ j\ge1$. Let \(k\ge 1\) and \(s>1\) and
\[
a:=\frac1{s-1},
\qquad
\psi(t):=
\begin{cases}
	0,& t\le 0,\\[1mm]
	\exp\!\bigl(-E_k(t^{-a})\bigr),& t>0.
\end{cases}
\]
We will prove that \(\psi\in G_k^s([-1,1])\) in Section \ref{sec9.1} and that $|m(\lambda_n)| \approx \log^{(k+2)} \lambda_n$ along a sequence $\lambda_n\to\infty$ in the following.

Since \(\psi(-t)=0\) for \(t>0\),
\[
m(\lambda)=\int_0^1 \frac{e^{i\lambda \psi(t)}-1}{t}\,dt.
\]
Set
\[
u=\psi(t)=\exp\!\bigl(-E_k(t^{-a})\bigr).
\]
Write \(x=t^{-a}\). Then
\[
\frac{du}{u}=-E_k'(x)\,dx
=
-\Bigl(\prod_{j=1}^k \log^{(j)}(1/u)\Bigr)\,dx
\]
and
\[
dx=-a\,x\,\frac{dt}{t}
=
-a\,\log^{(k+1)}(1/u)\,\frac{dt}{t}.
\]
Hence
\[
\frac{dt}{t}
=
(s-1)w(u)du,
\]
where $$w(u)=
\frac1{u\,\log(1/u)\,\log^{(2)}(1/u)\cdots \log^{(k+1)}(1/u)}.$$
Therefore
\[
m(\lambda)
=
(s-1)\int_0^{u_0}
(e^{i\lambda u}-1)w(u)
du
\]
for some fixed \(u_0\in(0,1)\).

Now take \(\lambda_n:=2\pi n\). Then
\[
-\Re m(\lambda_n)
=
(s-1)\int_0^{u_0}
(1-\cos(2\pi n u))w(u)du.
\]
Let
\[
I_j:=\Bigl[\frac{j+1/4}{n},\,\frac{j+3/4}{n}\Bigr],\qquad j=1,2,3,\dots.
\]
On \(I_j\), we have \(1-\cos(2\pi n u)\ge 1\). Since the weight $w(u)$ is decreasing for small \(u\), summing over \(1\le j\ls n\) yields for some fixed \(c_0\in(0,u_0)\)
\[
-\Re m(\lambda_n)
\gtrsim (s-1)
\int_{1/n}^{c_0} w(u)\,du\gs (s-1)\log^{(k+2)}n.
\]

Thus
\[
|m(\lambda_n)|
\ge -\Re m(\lambda_n)
\gtrsim (s-1)\log^{(k+2)}\lambda_n.
\]
This proves the sharpness of the upper bound $O(\log^{(k+2)}\la)$.

\section{Beyond Gevrey classes}\label{sec7}
In this section, we investigate classes of smooth functions that lie outside every Gevrey class. Jézéquel \cite{jez} proved a trace formula conjectured by Dyatlov--Zworski \cite{dz2016} for Anosov flows in dynamical systems that holds for certain intermediate regularity classes. See also \cite{CL,ptt2015,ptt2018}.

Let $I=[-1,1]$. Let $M=(M_n)_{n\ge 0}$ be a positive sequence with $M_n=\exp(cn^\alpha)$ ($\alpha>1$, $c>0$). It is log-convex. We consider the Denjoy--Carleman class $\mathcal{C}_M$ associated with this sequence. If $\psi\in \mathcal{C}_M(I)$, then there exists a constant $K>0$ such that for all $n\ge 0$
\begin{equation}\label{defdc1}
	\sup_{t\in I}|\psi^{(n)}(t)|\le K^{n+1}\exp(cn^\alpha).
\end{equation}
It is essentially the class used by Jézéquel \cite{jez}, and it is larger than any Gevrey class since $\exp(cn^\alpha)\gg (n!)^s$ for any $\alpha,s>1$ and $c>0$.

\begin{theorem}\label{thm:main4}
		Let $\psi\in \mathcal{C}_M(I)$. Then we have $m(\la)=O((\log\la)^{1-\frac1\alpha})$. 
	This bound is sharp: there exists a real-valued $\psi\in \mathcal{C}_M(I)$ such that
	\[
	|m(\lambda_n)| \approx (\log\la_n)^{1-\frac1\alpha}
	\quad\text{along a sequence }\lambda_n\to\infty.
	\]
\end{theorem}
The following flat-point estimate is the key to the theorem.
\begin{proposition}[Flat-point estimate]\label{prop4}
Let $\phi\in \mathcal{C}_M(I)$ be flat at $0$, i.e. $\phi^{(n)}(0)=0\ (n\ge 0).$
	Then there exist constants $A,B>0$ such that for all sufficiently small $|t|$,
	\begin{equation}\label{ub4}
		|\phi(t)|\le A\exp\Bigl(-B \Big(\log\frac1{|t|}\Big)^{\frac\alpha{\alpha-1}}\Bigr).
	\end{equation}
\end{proposition}
The estimate \eqref{ub4} is essentially sharp, since $\exp\Bigl(- \Big(\log\frac1{|t|}\Big)^{\frac\alpha{\alpha-1}}\Bigr)$ belongs to $\mathcal{C}_M(I)$ (see Section \ref{7sharp}). It can be proved by the Taylor--Legendre method in Remark \ref{rm3}. Indeed, by Taylor's theorem and minimization, we obtain for some constant $B>0$
\[|\phi(t)|\le \inf_{n\ge1}\frac{K^{n+1}\exp(cn^\alpha)|t|^n}{n!}\approx \exp\Bigl(-B \Big(\log\frac1{|t|}\Big)^{\frac\alpha{\alpha-1}}\Bigr),\quad |t|\ll1,\]
where the minimum is achieved at $n_*\approx (y/\alpha)^{\frac1{\alpha-1}}$ with $y=\log\frac1{K|t|}$. 

\begin{remark}\label{rm4}{\rm Applying Proposition \ref{prop:flat} gives only a weaker upper bound:
	\[|\phi(t)|\ls \exp\Bigl(-B \Big(\log\frac1{|t|}\Big)^{\frac1{\alpha-1}}\Bigr),\quad |t|\ll1.\]
	Roughly speaking, the Taylor--Legendre method works well for the Gevrey classes and other larger classes, while Proposition \ref{prop:flat} (Bang's Lemma \ref{lem:BB}) works well for the Gevrey classes and the refined Gevrey classes. See also Remark \ref{rm3} and Bang \cite[p.~143]{bang}. From this perspective, the Gevrey class is exactly the borderline between these two methods.}
\end{remark}

\subsection{Proof of the upper bound in Theorem \ref{thm:main4}}
By the reduction above, it suffices to consider the
case in which \(\phi\) is flat at \(0\). Let $\beta=\alpha/(\alpha-1)$. Then Proposition~\ref{prop4} gives
constants \(A,B>0\) and \(\delta\in(0,1)\) such that
\begin{equation}\label{eq:flat-small}
	|\phi(t)|\le A\exp\!\Bigl(-B\bigl(\log(1/t)\bigr)^{\beta}\Bigr)
	\qquad (0<t\le \delta).
\end{equation}
Using \(|e^{ix}-e^{iy}|\le \min\{2,|x-y|\}\), we obtain
\[
|m(\lambda)|
\ls 1+\int_0^\delta \min\{1,\lambda|\phi(t)|\}\,\frac{dt}{t}.
\]
Let $u_\la>0$ solve $\la\exp(-Bu^\beta)=1$. Then $u_\la\approx (\log\la)^{1/\beta}$. We split the interval and change variables to obtain
\begin{align*}
|m(\la)| 
&\ls u_\la+ \int_{u_\la}^\infty \la e^{-Bu^\beta}du\\
&\ls u_\la+\int_{u_\la}^\infty \la e^{-Bu_\la^{\beta-1}u}du\\
&\ls u_\la+u_\la^{1-\beta}\ls u_\la.
\end{align*}
This proves the upper bound in Theorem \ref{thm:main4}.


\subsection{Sharpness}\label{7sharp}  
Let $\alpha>1$ and $\beta=\alpha/(\alpha-1)$ and 
\[\psi(t)=\begin{cases}
	0,& t\le 0,\\[1mm]
	\exp(-(\log\frac1{t})^{\beta}),& 0<t\le1/e.
\end{cases}\]
We may take $I=[-1/e,1/e]$. It is harmless and simplifies the calculation.
We show that $\psi\in \mathcal{C}_M(I)$ in Section \ref{sec9.2}
 and that $|m(\lambda_n)| \approx (\log\la_n)^{1-\frac1\alpha}$
along a sequence $\lambda_n\to\infty$ in the following.

Since $\psi(-t)=0$ for $t>0$,
\[
m(\lambda)=\int_0^{e^{-1}}\frac{e^{i\lambda\psi(t)}-1}{t}\,dt.
\]
Set
\[
s:=\psi(t)=e^{-(\log(1/t))^\beta}.
\]

Differentiating gives
\[
\frac{dt}{t}
=
\frac1\beta \frac{ds}{s(\log(1/s))^{1-\frac1\beta}}
=
\frac1\beta \frac{ds}{s(\log(1/s))^{1/\alpha}}.
\]
Write
\[
w(s):=\frac{1}{s(\log(1/s))^{1/\alpha}},
\qquad 0<s\le e^{-1}.
\]


Take $\lambda_n:=2\pi n.$
Then
\[
-\Re m(\lambda_n)
=
\frac1\beta \int_0^{e^{-1}} (1-\cos(2\pi n s))\,w(s)\,ds.
\]
Fix
\[
c_0:=\frac1{4e},
\qquad
N_n:=\lfloor c_0 n\rfloor.
\]
For $1\le j\le N_n$ set
\[
I_{j,n}:=\Bigl[\frac{j+1/4}{n},\,\frac{j+3/4}{n}\Bigr].
\]
Since $(j+3/4)/n\le c_0+1/n<e^{-1}$ for all large $n$, these intervals lie inside $(0,e^{-1})$.
Moreover, on $I_{j,n}$ one has
\[
1-\cos(2\pi n s)\ge 1.
\]
Therefore
\[
-\Re m(\lambda_n)\ge \frac1\beta \sum_{j=1}^{N_n}\int_{I_{j,n}} w(s)\,ds.
\]
Because $w$ is decreasing,
\[
\int_{I_{j,n}} w(s)\,ds
\ge
\frac{1}{2n}w\Bigl(\frac{j+1}{n}\Bigr)
\ge
\frac12\int_{(j+1)/n}^{(j+2)/n} w(s)\,ds.
\]
Summing in $j$ gives
\[
-\Re m(\lambda_n)
\ge
\frac{1}{2\beta}\int_{2/n}^{(N_n+2)/n} w(s)\,ds
\gtrsim
\int_{2/n}^{c_0} \frac{ds}{s(\log(1/s))^{1/\alpha}}\approx
(\log n)^{1-\frac1\alpha}.
\]
Since $\lambda_n=2\pi n,$
\[
|m(\lambda_n)|\ge -\Re m(\lambda_n)\gtrsim (\log \la_n)^{1-\frac1\alpha}.
\]
This completes the proof.


\section{Further discussions}\label{sec8}
In this section, we consider the Denjoy--Carleman classes that are substantially larger than the Gevrey and intermediate classes discussed earlier. Let $I=[-1,1]$, and let $M=(M_n)_{n\ge0}$ be a positive sequence of the form $M_n=\exp^{(k)}(c n^\alpha)$, where $c,\alpha>0$ and $k\ge2$ is an integer. For large $n$, this sequence is log-convex, so it defines a Denjoy--Carleman class $\mathcal C_M$. If $\psi\in\mathcal C_M(I)$, then there exists a constant $K>0$ such that
\begin{equation}\label{defdc2}
    \sup_{t\in I} |\psi^{(n)}(t)|
    \le K^{n+1}\exp^{(k)}(c n^\alpha),
    \qquad n\ge0.
\end{equation}
Our goal is to understand how the oscillatory integral bound changes in this case. 
In what follows, we show that as the function classes become larger, the maximal growth of $m(\la)$ increases to the universal $C^\infty$ bound $o(\log\la)$.

\begin{theorem}\label{thm:main5}
	Let  $\psi\in \mathcal{C}_M(I)$. Then $m(\la)=O((\log\la)/(\log^{(k)}\la)^{\frac1\alpha})$. This bound is sharp: there exists a real-valued $\psi\in \mathcal{C}_M(I)$ such that
	\[
	|m(\lambda_n)| \approx (\log\la_n)/(\log^{(k)}\la_n)^{\frac1\alpha}
	\quad\text{along a sequence }\lambda_n\to\infty.
	\]

\end{theorem}
As before, the key ingredient is a flat-point estimate.
\begin{proposition}[Flat-point estimate]\label{prop5}
	Let $\phi\in \mathcal{C}_M(I)$ be flat at $0$, i.e. $\phi^{(n)}(0)=0\ (n\ge 0).$
	Then there exist constants $A,B>0$ such that for all sufficiently small $|t|$,
\begin{equation}\label{ub5}
	|\phi(t)|\le A\exp\Bigl(-B \Big(\log\frac1{|t|}\Big)\Big(\log^{(k)}\frac1{|t|}\Big)^{\frac1\alpha}\Bigr).
\end{equation}
\end{proposition}
The estimate \eqref{ub5} is essentially sharp, since $\exp\Bigl(- \Big(\log\frac1{|t|}\Big)\Big(\log^{(k)}\frac1{|t|}\Big)^{\frac1\alpha}\Bigr)$ belongs to $\mathcal{C}_M(I)$ (see Section \ref{8sharp}). It can be proved by the Taylor--Legendre method in Remark \ref{rm3}. Indeed, by Taylor's theorem and minimization, we obtain for some constant $B>0$
\[|\phi(t)|\le \inf_{n\ge1}\frac{K^{n+1}\exp^{(k)}(cn^\alpha)|t|^n}{n!}\approx \exp\Bigl(-B \Big(\log\frac1{|t|}\Big)\Big(\log^{(k)}\frac1{|t|}\Big)^{\frac1\alpha}\Bigr),\quad |t|\ll1,\]
where the minimum is achieved at $n_*\approx (\log^{(k-1)} y) ^{\frac1\alpha}$ with $y=\log\frac1{K|t|}$. 


\subsection{Proof of the upper bound in Theorem \ref{thm:main5}}
By the reduction above, it suffices to consider the case in which \(\phi\) is flat at \(0\). Then Proposition~\ref{prop5} gives constants \(A,B>0\) and \(\delta\in(0,1)\) such that
\begin{equation}
	|\phi(t)|\le A\exp\Bigl(-B \Big(\log\frac1{t}\Big)\Big(\log^{(k)}\frac1{t}\Big)^{\frac1\alpha}\Big)
	\qquad (0<t\le \delta).
\end{equation}
Using \(|e^{ix}-e^{iy}|\le \min\{2,|x-y|\}\), we obtain
\[
|m(\lambda)|
\ls 1+\int_0^\delta \min\{1,\lambda|\phi(t)|\}\,\frac{dt}{t}.
\]
Let $u_\la>0$ solve $\la\exp(-Bu(\log^{(k-1)}u)^{\frac1\alpha})=1$. Then $u_\la\approx \frac{\log\la}{(\log^{(k)}\la)^{\frac1\alpha}}$. 

We split the interval and change variables to obtain
\begin{align*}
	|m(\la)|
	&\ls u_\la+ \int_{u_\la}^\infty \la \exp(-B u(\log^{(k-1)}u)^{\frac1\alpha})du\\
	&\ls u_\la+\int_{u_\la}^\infty \la\exp(-B(\log^{(k-1)}u_\la)^{\frac1\alpha}u)du\\
	&\ls u_\la+(\log^{(k)}\la)^{-1/\alpha}\ls u_\la.
\end{align*}
This proves the upper bound in Theorem \ref{thm:main5}.

\subsection{Sharpness}\label{8sharp}
 For $x$ sufficiently large we write
\[
L_1(x):=\log x,
\qquad
L_{j+1}(x):=\log L_j(x)\quad (j\ge 1),
\]
and
\[
E_1(x):=e^x,
\qquad
E_{j+1}(x):=\exp(E_j(x))\quad (j\ge 1).
\]
Fix $k\ge 2$, $\alpha>0$, and $0<\delta\ll1$.
Let $I=(-\delta,\delta)$ and define
\[
\psi(t):=
\begin{cases}
	0, & -\delta<t\le 0,\\[4pt]
	\exp\!\Bigl(-\bigl(\log(1/t)\bigr)\,L_k(1/t)^{1/\alpha}\Bigr), & 0<t<\delta.
\end{cases}
\]
We will show that $\psi\in \mathcal{C}_M(I)$ in Section \ref{sec9.3} 
and that $|m(\lambda_n)| \approx (\log\la_n)/(\log^{(k)}\la_n)^{\frac1\alpha}$
along a sequence $\lambda_n\to\infty$ in the following.

Since $k\ge 2$, for $u:=\log(1/t)$ this may be rewritten as
\[
\psi(t)=e^{-Q(u)},
\qquad
Q(u):=uL_{k-1}(u)^{1/\alpha}.
\]


Since $\psi(-t)=0$ for $t>0$,
\[
m(\lambda)=\int_0^{\delta}\frac{e^{i\lambda\psi(t)}-1}{t}\,dt.
\]
Set
\[
s:=\psi(t)=e^{-Q(u)}\in (0,s_0],
\qquad
s_0:=\psi(\delta)>0.
\]
Hence
\begin{equation}
	\label{eq:m-omega}
	m(\lambda)=\int_0^{s_0}(e^{i\lambda s}-1)\,\omega(s)\,ds,
	\qquad
	\omega(s):=\frac{1}{sQ'(u(s))}.
\end{equation}


For $0<s\le s_0\ll1$, 
\begin{equation}
	\label{eq:weight-comparison}
\omega(s)\approx 
	\frac{1}{sL_k(1/s)^{1/\alpha}}.
\end{equation}



Take $\lambda_n:=2\pi n.$
By \eqref{eq:m-omega},
\[
-\Re m(\lambda_n)=\int_0^{s_0}(1-\cos(2\pi ns))\,\omega(s)\,ds.
\]
Using the lower bound in \eqref{eq:weight-comparison},
\[
-\Re m(\lambda_n)
\gtrsim
\int_0^{s_0}(1-\cos(2\pi ns))\,\frac{ds}{sL_k(1/s)^{1/\alpha}}.
\]
Choose a fixed number $0<c_*<s_0/4$ and let $N_n:=\lfloor c_*n\rfloor$.
For $1\le j\le N_n$ set
\[
I_{j,n}:=\Bigl[\frac{j+1/4}{n},\,\frac{j+3/4}{n}\Bigr].
\]
For large $n$, all these intervals lie in $(0,s_0]$, and on each $I_{j,n}$ one has
\[
1-\cos(2\pi ns)\ge 1.
\]
Therefore
\[
-\Re m(\lambda_n)
\gtrsim
\sum_{j=1}^{N_n}\int_{I_{j,n}}\frac{ds}{sL_k(1/s)^{1/\alpha}}.
\]
The weight $s\mapsto \bigl(sL_k(1/s)^{1/\alpha}\bigr)^{-1}$ is decreasing on $(0,s_0]$, so
\[
\int_{I_{j,n}}\frac{ds}{sL_k(1/s)^{1/\alpha}}
\ge \frac12\int_{(j+1)/n}^{(j+2)/n}\frac{ds}{sL_k(1/s)^{1/\alpha}}.
\]
Summing in $j$ gives
\[
-\Re m(\lambda_n)
\gtrsim
\int_{2/n}^{c_*}\frac{ds}{sL_k(1/s)^{1/\alpha}}\approx
\frac{\log n}{L_k(n)^{1/\alpha}}.
\]

Hence
\begin{equation}
	\label{eq:lower-seq}
	|m(\lambda_n)|\ge -\Re m(\lambda_n)
	\gtrsim
	\frac{\log \lambda_n}{L_k(\lambda_n)^{1/\alpha}}.
\end{equation}
\section{Derivative bounds}\label{sec9}


\subsection{Sharp example in Section \ref{6harp}}\label{sec9.1}
Let $E_1(x)=e^x,\ E_{j+1}(x)=\exp(E_j(x)),\ j\ge1$. Let \(k\ge 1\) and \(s>1\) and
\[
a:=\frac1{s-1},
\qquad
\psi(t):=
\begin{cases}
	0,& t\le 0,\\[1mm]
	\exp\!\bigl(-E_k(t^{-a})\bigr),& t>0.
\end{cases}
\]
We will prove that \(\psi\in G_k^s([-1,1])\).
For \(x\ge 1\), write
\[
M_k(x):=\prod_{j=1}^{k-1}E_j(x),
\qquad
\Omega_k(x):=x\,M_k(x)
\]
(with the convention \(M_1(x)=1\), so \(\Omega_1(x)=x\)).

\begin{lemma}
	For every \(k\ge 1\) there exists \(A_k\ge 1\) such that for all \(r\ge 0\) and all
	\(x\ge 1\),
	\[
	E_k^{(r)}(x)\le r!\,A_k^r\,M_k(x)^r\,E_k(x).
	\]
\end{lemma}

\begin{proof}
	We argue by induction on \(k\).
	
	For \(k=1\), \(E_1^{(r)}(x)=e^x=E_1(x)\), so the claim is immediate.
	
	Assume it holds for \(k-1\). Since \(E_k=e^{E_{k-1}}\), the Taylor expansion
	of \(E_k(x+h)\) at \(x\) gives
	\[
	\sum_{r=0}^\infty \frac{E_k^{(r)}(x)}{r!}h^r
	=
	E_k(x)\exp\!\bigl(E_{k-1}(x+h)-E_{k-1}(x)\bigr).
	\]
	Taking absolute values and using the induction hypothesis,
	\[
	\sum_{r=0}^\infty \frac{|E_k^{(r)}(x)|}{r!}|h|^r
	\le
	E_k(x)\exp\!\left(
	\sum_{m=1}^\infty A_{k-1}^m M_{k-1}(x)^m E_{k-1}(x)\,|h|^m
	\right).
	\]
	Since \(M_k(x)=E_{k-1}(x)M_{k-1}(x)\), choosing
	\[
	|h|=\frac1{2A_{k-1}M_k(x)}
	\]
	makes the exponent bounded by \(1\). Hence
	\[
	\sum_{r=0}^\infty \frac{|E_k^{(r)}(x)|}{r!}|h|^r \le e\,E_k(x).
	\]
	Comparing coefficients gives
	\[
	E_k^{(r)}(x)\le r!\,(2eA_{k-1})^r M_k(x)^r E_k(x).
	\]
	So the induction closes with \(A_k:=2eA_{k-1}\).
\end{proof}

Now put
\[
x(t):=t^{-a}\qquad (0<t\le 1),
\qquad
g(t):=E_k(x(t)),
\qquad
e^{-g(t)}=\psi(t)\quad (t>0).
\]
We first estimate derivatives of \(x(t)\). Since
\[
x^{(m)}(t)=(-1)^m a(a+1)\cdots (a+m-1)\,t^{-a-m},
\]
there is a constant \(B_a\ge 1\) such that for all \(m\ge 1\),
\[
|x^{(m)}(t)|\le m!\,B_a^m\,t^{-m}x(t).
\]

\begin{lemma}
	There exists \(C_1\ge 1\) such that for all \(n\ge 0\) and all \(0<t\le 1\),
	\[
	|g^{(n)}(t)|\le n!\,C_1^n\,t^{-n}\Omega_k(x(t))^n\,g(t).
	\]
\end{lemma}

\begin{proof}
	Using the Taylor expansion of \(E_k(x(t+h))\) around \(x(t)\),
	\[
	\sum_{n=0}^\infty \frac{g^{(n)}(t)}{n!}h^n
	=
	\sum_{r=0}^\infty \frac{E_k^{(r)}(x(t))}{r!}\,
	\bigl(x(t+h)-x(t)\bigr)^r .
	\]
	Taking absolute values and using the previous lemma,
	\[
	\sum_{n=0}^\infty \frac{|g^{(n)}(t)|}{n!}|h|^n
	\le
	g(t)\sum_{r=0}^\infty
	\left(
	A_k M_k(x(t))
	\sum_{m=1}^\infty \frac{|x^{(m)}(t)|}{m!}|h|^m
	\right)^r .
	\]
	Also,
	\[
	\sum_{m=1}^\infty \frac{|x^{(m)}(t)|}{m!}|h|^m
	\le
	x(t)\sum_{m=1}^\infty (B_a t^{-1}|h|)^m
	=
	x(t)\frac{B_a t^{-1}|h|}{1-B_a t^{-1}|h|}.
	\]
	Choose
	\[
	|h|=\frac1{4A_kB_a\,t^{-1}\Omega_k(x(t))}.
	\]
	Then \(B_a t^{-1}|h|\le 1/4\), hence the geometric factor above is at most
	\(2x(t)B_a t^{-1}|h|\), and therefore
	\[
	A_kM_k(x(t))
	\sum_{m=1}^\infty \frac{|x^{(m)}(t)|}{m!}|h|^m
	\le \frac12 .
	\]
	So
	\[
	\sum_{n=0}^\infty \frac{|g^{(n)}(t)|}{n!}|h|^n \le 2g(t).
	\]
	Comparing coefficients yields the claim.
\end{proof}

\begin{lemma}
	There exists \(C_2\ge 1\) such that for all \(n\ge 0\) and all \(0<t\le 1\),
	\[
	|\psi^{(n)}(t)|
	\le
	n!\,C_2^n\,t^{-n}\Omega_k(x(t))^n\,e^{-g(t)/2}.
	\]
\end{lemma}

\begin{proof}
	Since \(\psi=e^{-g}\),
	\[
	\sum_{n=0}^\infty \frac{\psi^{(n)}(t)}{n!}h^n
	=
	e^{-g(t)}\exp\!\bigl(-(g(t+h)-g(t))\bigr).
	\]
	Taking absolute values and using the previous lemma,
	\[
	\sum_{n=0}^\infty \frac{|\psi^{(n)}(t)|}{n!}|h|^n
	\le
	e^{-g(t)}
	\exp\!\left(
	\sum_{m=1}^\infty C_1^m t^{-m}\Omega_k(x(t))^m g(t)\,|h|^m
	\right).
	\]
	Choose
	\[
	|h|=\frac1{4C_1\,t^{-1}\Omega_k(x(t))}.
	\]
	Then the series in the exponent is at most \(g(t)/2\), so
	\[
	\sum_{n=0}^\infty \frac{|\psi^{(n)}(t)|}{n!}|h|^n
	\le e^{-g(t)/2}.
	\]
	Comparing coefficients gives
	\[
	|\psi^{(n)}(t)|
	\le
	n!\,(4C_1)^n\,t^{-n}\Omega_k(x(t))^n\,e^{-g(t)/2}.
	\]
\end{proof}

Because \(a=1/(s-1)\), we have \(t^{-1}=x(t)^{\,s-1}\). Hence
\[
t^{-1}\Omega_k(x(t))
=
x(t)^s \prod_{j=1}^{k-1}E_j(x(t)).
\]
Therefore
\begin{equation}\label{eq:star}
	|\psi^{(n)}(t)|
	\le
	n!\,C_2^n
	\left(
	x(t)^s\prod_{j=1}^{k-1}E_j(x(t))
	\right)^n
	e^{-E_k(x(t))/2}.
\end{equation}
Now choose $n_0$ large and define
\[
x_n:=\log^{(k)}\bigl((n+n_0)^2\bigr).
\]
Then
\[
E_k(x_n)=(n+n_0)^2,
\qquad
x_n\approx \log^{(k)}n,
\qquad
E_j(x_n)=\log^{(k-j)}\bigl((n+n_0)^2\bigr)\approx \log^{(k-j)}n
\]
for $1\le j\le k-1$. Hence
\[
x_n^s\prod_{j=1}^{k-1}E_j(x_n)\lesssim Q_{k,s}(n).
\]

If $1\le x\le x_n$, then by monotonicity,
\[
x^s\prod_{j=1}^{k-1}E_j(x)\le x_n^s\prod_{j=1}^{k-1}E_j(x_n)\lesssim Q_{k,s}(n).
\]

If $x\ge x_n$, set
\[
R(x):=x^s\prod_{j=1}^{k-1}E_j(x).
\]
Since
\[
\log R(x)=s\log x+x+E_1(x)+\cdots+E_{k-2}(x)=o(E_{k-1}(x)),
\]
we have $R(x)\le E_k(x)^{1/2}$ for all large $x$. Therefore, writing $y=E_k(x)$,
\[
R(x)^n e^{-E_k(x)/2}\le y^{n/2}e^{-y/2}.
\]
As $y\ge E_k(x_n)=(n+n_0)^2\ge n$, the function $y^{n/2}e^{-y/2}$ is decreasing, so
\[
y^{n/2}e^{-y/2}\le (n+n_0)^n e^{-(n+n_0)^2/2}\ls1.
\]
Combining this with \eqref{eq:star}, we obtain
\[
\sup_{0<t\le 1}|\psi^{(n)}(t)|\le K^{n+1} n!Q_{k,s}(n)^n
\]
for some constant $K>0$ and all large $n$. Since \(\psi^{(n)}(t)\to 0\) as
\(t\downarrow 0\) by \eqref{eq:star}, the extension by \(0\) to \((-\infty,0]\) is
\(C^\infty\) and flat at \(0\). Therefore \(\psi\in G_k^s([-1,1])\).


\subsection{Sharp example in Section \ref{7sharp}}\label{sec9.2}
Let $\alpha>1$ and $\beta=\alpha/(\alpha-1)$ and 
\[\psi(t)=\begin{cases}
	0,& t\le 0,\\[1mm]
	\exp(-(\log\frac1{t})^{\beta}),& 0<t\le1/e.
\end{cases}\]
We take $I=[-1/e,1/e]$ and show that $\psi\in \mathcal{C}_M(I)$ with $M_n=\exp(c n^\alpha)$ for some $c>0$.

For $0<t<e^{-1}$ set
\[
u:=\log(1/t)\ge 1,
\qquad
g(u):=e^{-u^\beta},
\]
so that $\psi(t)=g(u)$.

\begin{lemma}
	\label{lem:dtdu}
	For every $n\ge 1$ there exist real numbers $a_{n,k}$, $1\le k\le n$, such that
	\[
	\frac{d^n}{dt^n}g(\log(1/t))
	=
	t^{-n}\sum_{k=1}^n a_{n,k}\,g^{(k)}(u),
	\]
	and
	\[
	\sum_{k=1}^n |a_{n,k}|\le n!.
	\]
\end{lemma}

\begin{proof}
	For $n=1$ this is immediate:
	\[
	\frac{d}{dt}g(\log(1/t))=-t^{-1}g'(u),
	\]
	so we may take $a_{1,1}=-1$.
	
	Assume the statement true for some $n\ge 1$.
	Differentiate
	\[
	\frac{d^n}{dt^n}g(\log(1/t))
	=
	t^{-n}\sum_{k=1}^n a_{n,k}\,g^{(k)}(u).
	\]
	Since $u'=-t^{-1}$, we get
	\[
	\frac{d}{dt}\Bigl(t^{-n}F(u)\Bigr)
	=
	-t^{-n-1}\bigl(nF(u)+F'(u)\bigr).
	\]
	Applying this with $F(u)=\sum_{k=1}^n a_{n,k}g^{(k)}(u)$ yields
	\[
	\frac{d^{n+1}}{dt^{n+1}}g(\log(1/t))
	=
	t^{-n-1}\sum_{k=1}^{n+1} a_{n+1,k}\,g^{(k)}(u),
	\]
	where, with the convention $a_{n,0}=a_{n,n+1}=0$,
	\[
	a_{n+1,k}=-(n a_{n,k}+a_{n,k-1}).
	\]
	Hence
	\[
	\sum_{k=1}^{n+1}|a_{n+1,k}|
	\le
	n\sum_{k=1}^n|a_{n,k}|+\sum_{k=1}^n|a_{n,k}|
	\le
	(n+1)n!=(n+1)!.
	\]
	This closes the induction.
\end{proof}

We next estimate $g^{(k)}$.

\begin{lemma}
	\label{lem:gk}
	There exists a constant $C_0\ge 1$ such that for all integers $k\ge 1$ and all $u\ge 1$,
	\[
	|g^{(k)}(u)|
	\le
	C_0^k\,k!\,k^k\,u^{(\beta-1)k}e^{-u^\beta}.
	\]
\end{lemma}

\begin{proof}
	Write
	\[
	h(u):=-u^\beta,
	\qquad
	g(u)=e^{h(u)}.
	\]
	For $m\ge 1$,
	\[
	h^{(m)}(u)=-(\beta)_m\,u^{\beta-m},
	\]
	where $(\beta)_m=\beta(\beta-1)\cdots (\beta-m+1)$.
	Since
	\[
	|\,\beta-j\,|\le |\beta|+j\le (1+|\beta|)(j+1)
	\qquad (j\ge 0),
	\]
	we obtain
	\[
	|(\beta)_m|
	\le (1+|\beta|)^m m!.
	\]
	Therefore, after enlarging the constant if needed,
	\begin{equation}
		\label{eq:hm}
		|h^{(m)}(u)|\le C_0^m m!\,u^{\beta-m}
		\qquad (m\ge 1,\ u\ge 1).
	\end{equation}
	
	Now apply Fa\`a di Bruno in the partition form:
	\[
	g^{(k)}(u)
	=
	e^{h(u)}
	\sum_{\pi\in \Pi_k}\prod_{B\in\pi} h^{(|B|)}(u),
	\]
	where $\Pi_k$ is the set of partitions of $\{1,\dots,k\}$.
	Fix $\pi\in\Pi_k$, let $r:=|\pi|$ (the number of blocks in $\pi$) and let $m_1,\dots,m_r$ be its block sizes, so
	$m_1+\cdots+m_r=k$.
	Using \eqref{eq:hm},
	\[
	\prod_{\ell=1}^r |h^{(m_\ell)}(u)|
	\le
	C_0^k\Bigl(\prod_{\ell=1}^r m_\ell!\Bigr)\,u^{r\beta-k}.
	\]
	Since $r\le k$ and $\prod_{\ell=1}^r m_\ell!\le k!$, we obtain
	\[
	\prod_{\ell=1}^r |h^{(m_\ell)}(u)|
	\le
	C_0^k k!\,u^{(\beta-1)k}.
	\]
	The number of partitions of a $k$-element set is the Bell number $B_k$, and $B_k\le k^k$.
	Hence
	\[
	|g^{(k)}(u)|
	\le
	C_0^k\,k!\,k^k\,u^{(\beta-1)k}e^{-u^\beta}.
	\]
\end{proof}

Combining Lemmas~\ref{lem:dtdu} and \ref{lem:gk}, for $n\ge 1$ and $0<t<e^{-1}$ we get
\[
|\psi^{(n)}(t)|
\le
t^{-n}\sum_{k=1}^n |a_{n,k}|\,|g^{(k)}(u)|
\le
t^{-n} e^{-u^\beta}\Bigl(\sum_{k=1}^n |a_{n,k}|\Bigr)
\max_{1\le k\le n}\bigl(C_0^k k!\,k^k\,u^{(\beta-1)k}\bigr).
\]
By Lemma~\ref{lem:dtdu},
\[
\sum_{k=1}^n|a_{n,k}|\le n!,
\]
and since $k\le n$ and $u\ge 1$,
\[
C_0^k k!\,k^k\,u^{(\beta-1)k}
\le
C_0^n n!\,n^n\,u^{(\beta-1)n}.
\]
Therefore
\[
|\psi^{(n)}(t)|
\le
C_0^n (n!)^2 n^n\,e^{nu}\,u^{(\beta-1)n}e^{-u^\beta}.
\]
Taking logarithms and using Stirling in the crude form $\log(n!)\lesssim n\log(n+1)$,
we obtain
\begin{equation}
	\label{eq:rough-bound}
	|\psi^{(n)}(t)|
	\le
	\exp\!\Bigl(C_1 n\log(n+1)+nu+(\beta-1)n\log u-u^\beta\Bigr)
\end{equation}
for a constant $C_1>0$.

For fixed $n$, the exponent on the right-hand side tends to $-\infty$ as $u\to\infty$.
Hence $\psi^{(n)}(t)\to 0$ as $t\downarrow 0$ for every $n$, so the extension by $0$ to
$t\le 0$ is $C^\infty$ at the origin.

It remains to optimize \eqref{eq:rough-bound}. Since $u\ge 1$, we have $\log u\le u$, hence
\[
nu+(\beta-1)n\log u \le \beta n u.
\]
Because $\alpha$ and $\beta$ are conjugate exponents,
Young's inequality gives, for every $\varepsilon\in(0,1)$,
\[
\beta n u \le \varepsilon u^\beta + C_\varepsilon n^\alpha.
\]
Choosing $\varepsilon=\frac12$,
\[
nu+(\beta-1)n\log u-u^\beta \le -\frac12 u^\beta + C_2 n^\alpha \le C_2 n^\alpha.
\]
Thus \eqref{eq:rough-bound} implies
\[
|\psi^{(n)}(t)|
\le
\exp\!\bigl(C_1 n\log(n+1)+C_2 n^\alpha\bigr)
\qquad (0<t<e^{-1}).
\]
Since $\alpha>1$, we have $n\log(n+1)\le \varepsilon n^\alpha + C_\varepsilon n$.
Absorbing the linear term into $K^{n+1}$, we obtain
\[
\sup_{t\in I}|\psi^{(n)}(t)|\le K^{n+1}e^{c n^\alpha}
\]
for suitable constants $K,c>0$ and all $n\ge 0$. This proves $\psi\in \mathcal{C}_M(I)$ in Section \ref{sec7}.


\subsection{Sharp example in Section \ref{8sharp}}\label{sec9.3}
For $x$ sufficiently large we write
\[
L_1(x):=\log x,
\qquad
L_{j+1}(x):=\log L_j(x)\quad (j\ge 1),
\]
and
\[
E_1(x):=e^x,
\qquad
E_{j+1}(x):=\exp(E_j(x))\quad (j\ge 1).
\]
Fix $k\ge 2$, $\alpha>0$, and choose $\delta>0$ so small that
\[
L_j(1/t)\ge 2
\qquad (0<t\le \delta,\ 1\le j\le k).
\]
Let $I:=(-\delta,\delta)$ and define
\[
\psi(t):=
\begin{cases}
	0, & -\delta<t\le 0,\\[4pt]
	\exp\!\Bigl(-\bigl(\log(1/t)\bigr)\,L_k(1/t)^{1/\alpha}\Bigr), & 0<t<\delta.
\end{cases}
\]
Since $k\ge 2$, for $u:=\log(1/t)$ this may be rewritten as
\[
\psi(t)=e^{-Q(u)},
\qquad
Q(u):=uA(u),
\qquad
A(u):=L_{k-1}(u)^{1/\alpha}.
\]
For $0<t<\delta$ set $g(u):=e^{-Q(u)}$.
To estimate $g^{(r)}$ we first record a direct bound on the derivatives of $A$ and $Q$.
Introduce the auxiliary functions
\[
\eta_0(u):=u^{-1},
\qquad
\eta_j(u):=L_j(u)^{-1}\quad (1\le j\le k-1).
\]
Observe that
\[
\eta_0'(u)=-\eta_0(u)^2,
\qquad
\eta_j'(u)=-\eta_0(u)\eta_1(u)\cdots \eta_{j-1}(u)\eta_j(u)^2
\quad (j\ge 1).
\]
In particular, each derivative introduces at least one extra factor of $\eta_0(u)=u^{-1}$.

\begin{lemma}
	\label{lem:AQ}
	There exists a constant $C\ge 1$ such that, for every integer $r\ge 1$ and every sufficiently large $u$,
	\[
	|A^{(r)}(u)|\le C^r r!\,A(u)u^{-r},
	\qquad
	|Q^{(r)}(u)|\le C^r r!\,A(u)u^{1-r}.
	\]
	Moreover,
	\[
	Q'(u)=A(u)\Bigl(1+\frac{1}{\alpha L_1(u)\cdots L_{k-2}(u)L_{k-1}(u)}\Bigr),
	\]
	so in particular
	\[
	Q'(u)\approx A(u)
	\qquad (u\to\infty).
	\]
\end{lemma}

\begin{proof}
	Since
	\[
	A(u)=L_{k-1}(u)^{1/\alpha}=\eta_{k-1}(u)^{-1/\alpha},
	\]
	a straightforward induction shows that for every $r\ge 0$,
	\[
	A^{(r)}(u)=A(u)u^{-r}P_r\bigl(\eta_1(u),\dots,\eta_{k-1}(u)\bigr),
	\]
	where $P_r$ is a polynomial with real coefficients and where the $\ell^1$-norm of the coefficient vector of $P_r$ is bounded by $C^rr!$ for a suitable constant $C$ independent of $r$.
	Indeed, differentiating the identity above produces three kinds of terms: one from differentiating $A$, one from differentiating $u^{-r}$, and one from differentiating the polynomial in the variables $\eta_j$.
	Each differentiation contributes one extra factor of $u^{-1}=\eta_0$, and the coefficient growth is at most linear in $r$ at each step; hence the total coefficient growth is bounded by $C^rr!$.
	Since $0<\eta_j(u)\le 1$ for large $u$, this gives
	\[
	|A^{(r)}(u)|\le C^rr!\,A(u)u^{-r}.
	\]
	
	Now $Q(u)=uA(u)$, so Leibniz' rule gives, for $r\ge 1$,
	\[
	Q^{(r)}(u)=uA^{(r)}(u)+rA^{(r-1)}(u).
	\]
	Using the bound just proved,
	\[
	|Q^{(r)}(u)|
	\le C^rr!A(u)u^{1-r}+rC^{r-1}(r-1)!A(u)u^{1-r}
	\le (2C)^rr!A(u)u^{1-r}.
	\]
	This proves the stated estimate after enlarging $C$.
	Finally,
	\[
	A'(u)=\frac{A(u)}{\alpha uL_1(u)\cdots L_{k-2}(u)L_{k-1}(u)},
	\]
	so
	\[
	Q'(u)=A(u)+uA'(u)
	=A(u)\Bigl(1+\frac{1}{\alpha L_1(u)\cdots L_{k-2}(u)L_{k-1}(u)}\Bigr),
	\]
	which implies $Q'(u)\approx A(u)$ for large $u$.
\end{proof}

\begin{lemma}
	\label{lem:g}
	There exists a constant $C\ge 1$ such that for every integer $r\ge 1$ and every sufficiently large $u$,
	\[
	|g^{(r)}(u)|\le C^r (r!)^2 A(u)^r e^{-Q(u)}.
	\]
\end{lemma}

\begin{proof}
	Apply Fa\`a di Bruno in the partition form:
	\[
	g^{(r)}(u)=e^{-Q(u)}\sum_{\pi\in\Pi_r}\prod_{B\in\pi}(-Q^{(|B|)}(u)),
	\]
	where $\Pi_r$ is the set of partitions of $\{1,\dots,r\}$.
	For a partition $\pi\in\Pi_r$, let $m_B:=|B|$.
	By Lemma~\ref{lem:AQ},
	\[
	\prod_{B\in\pi}|Q^{(m_B)}(u)|
	\le C^r\Bigl(\prod_{B\in\pi}m_B!\Bigr)A(u)^{|\pi|}
	\le C^r r!\,A(u)^r,
	\]
	because $|\pi|\le r$ and $\prod_B m_B!\le r!$.
	Summing over partitions and using the Bell-number bound $|\Pi_r|\le e^r r!$, we obtain
	\[
	|g^{(r)}(u)|\le (Ce)^r(r!)^2A(u)^re^{-Q(u)}.
	\]
	Absorb $e$ into the constant.
\end{proof}

We may now estimate $\psi^{(n)}$ directly.
By Lemmas~\ref{lem:dtdu} and \ref{lem:g}, for $n\ge 1$ and $0<t<\delta$,
\[
|\psi^{(n)}(t)|
\le t^{-n}\sum_{r=1}^n |a_{n,r}|\,|g^{(r)}(u)|
\le t^{-n}n!\max_{1\le r\le n}|g^{(r)}(u)|.
\]
Hence
\[
|\psi^{(n)}(t)|
\le C^n (n!)^3 e^{nu}A(u)^n e^{-uA(u)}.
\]
Taking logarithms and using $\log(n!)\le C_0n\log(n+1)$, we get
\begin{equation}
	\label{eq:basic-exponent}
	|\psi^{(n)}(t)|
	\le \exp\!\Bigl(C_1n\log(n+1)+nu+n\log A(u)-uA(u)\Bigr).
\end{equation}
We now optimize the exponent.

Set
\[
F_n(u):=nu+n\log A(u)-uA(u).
\]
There are two cases.

\smallskip
\emph{Case 1: $A(u)\le 2n$.}
Then $L_{k-1}(u)=A(u)^\alpha\le (2n)^\alpha$, hence
\[
u\le E_{k-1}((2n)^\alpha).
\]
Therefore
\[
F_n(u)\le nE_{k-1}((2n)^\alpha)+n\log(2n).
\]
Since $k\ge 2$, the function $E_{k-1}(x)$ dominates polynomial factors; thus, after enlarging the constant,
\[
F_n(u)\le E_{k-1}(C_2n^\alpha).
\]

\smallskip
\emph{Case 2: $A(u)\ge 2n$.}
Since $\log A(u)=\frac1\alpha L_k(u)=o(u)$, for large $u$,
$\log A(u)\le \frac{u}{2}$. If $A(u)\ge 2n$, then
\[
F_n(u)
=nu+n\log A(u)-uA(u)
\le nu+\frac{n}{2}u-uA(u)
\le \frac32 nu-2nu
= -\frac12 nu\le 0.
\]

% Since $A(u)=L_{k-1}(u)^{1/\alpha}\le u^{1/\alpha}$, this implies $u\ge (2n)^\alpha\ge n$.
% Also, for large $u$ we have $A(u)\ge 2$, so $\log A(u)\le A(u)/2$.
% Hence
% \[
% F_n(u)
% =-u(A(u)-n)+n\log A(u)
% \le -\frac12uA(u)+n\log A(u)
% \le -\frac12nA(u)+n\log A(u)
% \le 0
% \]
%for all sufficiently large $u$.
Thus in this case as well,
\[
F_n(u)\le E_{k-1}(C_2n^\alpha).
\]

Substituting into \eqref{eq:basic-exponent}, we obtain
\[
|\psi^{(n)}(t)|\le \exp\!\bigl(C_1n\log(n+1)+E_{k-1}(C_2n^\alpha)\bigr).
\]
Because $k\ge 2$, one has $n\log(n+1)\le E_{k-1}(C_3n^\alpha)$ for a larger constant $C_3$, and therefore
\[
|\psi^{(n)}(t)|\le \exp\!\bigl(E_{k-1}(cn^\alpha)\bigr)=E_k(cn^\alpha)
\]
for some $c>0$.
After increasing the constants to account for finitely many small values of $n$, we obtain
\[
\sup_{0<t<\delta}|\psi^{(n)}(t)|\le K^{n+1}E_k(cn^\alpha).
\]
Since the exponent in \eqref{eq:basic-exponent} tends to $-\infty$ as $u\to\infty$ for each fixed $n$, all derivatives tend to $0$ as $t\downarrow 0$; thus the extension by $0$ to $t\le 0$ is $C^\infty$ at the origin. This proves $\psi\in \mathcal{C}_M(I)$ in Section \ref{sec8}.

\bibliographystyle{plain}
\begin{thebibliography}{0}
	\bibitem{acp14}H. M. Al-Qassem, L. Cheng, and Y. Pan. Rough oscillatory singular integrals on $\R^n$. Studia Math. 221 (2014), no. 3, 249-267.
	\bibitem{acp25}H. M. Al-Qassem, L. Cheng and Y. Pan. Uniform boundedness of oscillatory singular integrals with rational phases.
	Forum Math. 37 (2025), no. 6, 1665-1672.
	\bibitem{bang}T. Bang, On quasi-analytic functions, Math. Scand., 1 (1953), 137-152.
	\bibitem{bruna}J. Bruna, An extension theorem of Whitney type for non quasi-analytic classes of functions, J. London Math. Soc. (2) 22 (1980), no. 3, 495-505.
	
	\bibitem{cp99}A. Carbery and S. P\'erez. Maximal functions and hilbert transforms along variable flat curves. Math. Res. Lett., 6:237-250, 1999.
	\bibitem{CL}M. Cicognani and  D. Lorenz. Strictly hyperbolic equations with coefficients low-regular in time and smooth in space. J. Pseudo-Differ.
	Oper. Appl. 2018, 9, 643-675.
	\bibitem{dz2016}S. Dyatlov and M. Zworski.  Dynamical zeta functions for Anosov flows via microlocal analysis. Ann. Sci.
	Éc. Norm. Supér. (4) 49(3), 543-577 (2016)
	\bibitem{fw03}M. Folch-Gabayet and J. Wright. An oscillatory integral estimate associated to rational phases.
	J. Geom. Anal. 13 (2003), no. 2, 291-299.
	\bibitem{fw032}M. Folch-Gabayet and J. Wright, An estimation for a family of oscillatory integrals, Studia Math. 154 (2003), no. 1, 89-97. 
	\bibitem{gev}M. Gevrey, Sur la nature analytique des solutions des équations aux dérivées partielles. Premier mémoire.
	Ann. Sci. École Norm. Sup. (3) 35 (1918), 129-190.
	\bibitem{hor1}L. H\"ormander. The Analysis of Linear Partial Differential Operators I Distribution Theory and
	Fourier Analysis, Springer-Verlag, Berlin, 1985.

	\bibitem{jez}M. Jézéquel. Global trace formula for ultra-differentiable Anosov flows. Commun. Math. Phys. 2021, 385, 1771-1834.
	\bibitem{nvww}A. Nagel, J. Vance, S. Wainger, and D. Weinberg. Hilbert transforms for convex curves. Duke Math.
	J, 50(3):735-744, 1983.
	\bibitem{nw}A. Nagel and S. Wainger, Hilbert transforms associated with plane curves, Trans. Amer. Math. Soc. 223 (1976), 235--252.
	\bibitem{pan91}Y. Pan. Uniform estimates for oscillatory integral operators. J. Funct. Anal.,
	100(1):207-220, 1991.
	\bibitem{pan93}Y. Pan. $L^2$ estimates for convolution operators with oscillating kernels. Math. Proc. Cambridge Philos. Soc., 113, 1993.
	\bibitem{pp10}M. Papadimitrakis and I. R. Parissis. Singular oscillatory integrals on $\mathbb{R}^n$.
	Math. Z. 266 (2010), no. 1, 169-179.
	\bibitem{pa08}I. R. Parissis. A sharp bound for the Stein-Wainger oscillatory integral. Proc. Amer.
	Math. Soc., pages 963-972, 2008.
	\bibitem{ps86}D.H. Phong and E. M. Stein. Hilbert integrals, singular integrals, and radon transforms I. Acta Math., 157(1):99-157, 1986.
	\bibitem{ptt2015}S.~Pilipović, N.~Teofanov, and F.~Tomić. On a class of ultradifferentiable functions. Novi Sad J. Math., 45:125--142, 2015.
	\bibitem{ptt2018}S.~Pilipović, N.~Teofanov, and F.~Tomić. Beyond Gevrey regularity: Superposition and propagation of
	singularities. Filomat, 32(8):2763--2782, 2018.
	\bibitem{rod93}L. Rodino. Linear Partial Differential Operators in Gevrey Spaces. World Scientific Publishing Co., Inc., River Edge, NJ, 1993.
	\bibitem{rs87}F. Ricci and E.M. Stein. Harmonic analysis on nilpotent groups and singular integrals I. oscillatory
	integrals. J. Funct. Anal., 73(1):179-194, 1987.
	\bibitem{seeger94}A. Seeger. $L^2$ estimates for a class of singular oscillatory integrals. Math. Res. Lett, 1:65-73, 1994
%	\bibitem{steinbook}E. M. Stein, Oscillatory integrals in Fourier analysis, Beijing lectures in harmonic analysis
%	(Beijing, 1984), Ann. of Math. Stud., vol. 112, Princeton Univ. Press, Princeton, NJ, 1986,	pp. 307–355.
	\bibitem{sw70}E. Stein and S. Wainger, The estimation of an integral arising in multiplier transformations, Studia Math. 35 (1970), 101-104.
	\bibitem{sw78}E. M. Stein and S. Wainger, Problems in harmonic analysis related to curvature, Bull. Amer. Math. Soc. 84 (1978), no. 6, 1239-1295. 
	\bibitem{ww23}C. Wang and H. Wu. A note on singular oscillatory integrals with certain rational phases.
	C. R. Math. Acad. Sci. Paris 361 (2023), 363-370.
	\bibitem{wz21}X. Wang and C. Zhang, Sharp endpoint estimates for eigenfunctions restricted to submanifolds of codimension 2. Adv. Math. 386 (2021), Paper No. 107835, 20 pp.
\end{thebibliography}
\end{document}
